\documentclass[12pt]{article}

\usepackage{hyperref}
\usepackage{natbib}
\bibpunct[: ]{(}{)}{;}{a}{}{,}
\usepackage[margin=1in]{geometry}
\usepackage{setspace} 
\usepackage{graphicx}
\usepackage{lscape}
\usepackage{subfigure}

\hypersetup {
     pdfauthor={Omar Lizardo},
     pdfpagemode = UseNone,
     pdfstartview=FitH,
     citebordercolor = 1 0 0,
        }
             
\begin{document}

\title{Cultural Capital in Interaction\footnote{Direct correspondence to Omar Lizardo, University of Notre Dame, Department of Sociology, 810 Flanner Hall, Notre Dame, IN, 46556  Tel: (574) 631-1855, Fax: (574) 631-9238, Email: \href{mailto:olizardo@nd.edu}{\tt olizardo@nd.edu}.}}

\author{Omar Lizardo\thanks{I would like to thank Gene Halton, Robert Fishman and Mary Ellen Konieczny for helpful comments on earlier drafts.}
\\\textit{University of Notre Dame}}

\date{\today}
\maketitle

\begin{abstract}
In this paper I analyze survey data from a probability sample of Americans to  for the first time document quantitatively the distribution and predictors of what I call ``culture-mediated interaction'': essentially, the use of cultural content to generate conversation with friends.  I argue that this phenomenon--central to many {\em theories} of the connection between culture and stratification--deserves more {\em empirical} attention. I show that culture-mediated interaction is a pervasive feature of the lives of many Americans, especially those who are culturally active. In addition, I show that rates of culture-mediated interaction appear to be stratified by the usual markers of status (such as educational attainment), and for some cultural domains, this stratification is independent of frequency of cultural engagement.  I close by outlining the theoretical and empirical implications of taking the phenomenon of culture-mediated interaction seriously. 
\end{abstract}
\onehalfspacing

\section{Introduction}

Contemporary empirical research at the intersection of the sociology of culture and stratification is mainly concerned  with exploring inequalities linked to social position in the relative rates of access to and engagement with symbolic goods  \citep{dimaggio00, vaneijck00, katzgerro04}.  The central issue  remains the extent to which taste and culture consumption oriented lifestyles continue to be associated with social position \citep{tomlinson03}.  This common concern can be appreciated whether the construct of ``social position'' is conceptualized in terms of educational attainment \citep{bourdieu_darbel91, bryson96, vaneijck99}, membership in broadly defined occupational groups \citep{bourdieu84, dimaggio_useem78a, dimaggio_useem78b, peterson_simkus93, tomlinson03}, aggregate measures of social class \citep{katzgerro_shavit98, katzgerro02, katzgerro06}, or constructed scalar indices of social status based on social linkages among and within occupational groups \citep{chan_goldthorpe05, chan_goldthorpe07a, chan_goldthorpe07b}. 

The overall conclusion reached in this still vibrant line of research is, that contrary to the predictions of ``postmodernization'' theories that foresaw the decline in the linkage between class and lifestyle in market-dominated ``consumer societies'' lifestyle practices continue to be strongly shaped---albeit in increasingly complex ways \citep{peterson05}---by social position \citep{holt98, katzgerro99, tomlinson03, tampubolon08, warde_etal08}.  Contemporary forms of cultural stratification continue to be structured according to markers associated with social status and education rather than economic class \citep{alderson_etal07,chan_goldthorpe07a, chan_goldthorpe07b, katzgerro02, katzgerro04, sullivan_katzgerro07, silva06, bennett_etal99}.  Confirming the basic findings of an earlier generation of research \citep{baumol_bowen66, dimaggio_useem78a, ganzeboom82}.

The strong focus of the cultural stratification literature on the question of {\em whether} lifestyle practices continue to be unevenly partitioned according to class and status markers, has resulted in a relative lack of attention to an equally important question:  what is the primary {\em role} (or ``function'') of aesthetic consumption in the everyday lives of members of different social strata?  More precisely, given that some groups---usually those who enjoy markers of privilege and status---are more likely to engage culture and the arts more than other groups, to what {\em uses} do they put the culture that the consume?  The relative neglect of the behavioral consequences of aesthetic consumption  appears even more surprising in light of the fact that the issue of the role played by engagement with symbolic goods was an continues to be central to the theoretical corpus from which most analysts of the culture-stratification linkage depart in the first place \citep{collins75, dimaggio87, douglas_isherwood96, fine77}.


A key insight to be gained from this (largely theoretical) tradition is that if we cannot get a solid theoretical handle on the uses to which persons put the cultural repertoires acquired through aesthetic experiences it is difficult to analytically justify their status as ``cultural capital'' in a given social setting \citep{mohr_dimaggio95, dimaggio96, dimaggio00, carter03}.  This is the reason why in their influential paper Lamont and Lareau \citeyearpar{lamont_lareau88} linked the notion of cultural capital to an explicit conceptual framework that highlighted the active deployment of a given repertoire of cultural tastes and lifestyle practices for the purposes of drawing symbolic boundaries. In a similar way, in key theoretical statements DiMaggio \citeyearpar{dimaggio87, dimaggio00} is clear---building on the early work of Collins \citeyearpar{collins75}---in connecting the aesthetic consumption of cultural goods to routine exploitation of cultural knowledge in social interaction within everyday settings (see also \citet{erickson96} and \citet{lizardo06a} for empirical applications of this idea).  It is thus safe to say that the set of cultural aptitudes that result in the relatively higher likelihood to consume cultural goods only becomes actualized as cultural capital when it becomes linked to everyday processes of informal socialization and thus status group formation in the Weberian sense \citep{lizardo06a}.

In this paper, I draw out the empirical implications of this set of starting ideas for our understanding of the {\em social uses} of aesthetic consumption in interaction.  My analytical departing point is that the everyday use of cultural knowledge and experiences gained through aesthetic consumption can be thought of as a social resource that becomes the primary fodder for social interaction in conversation in both relatively formal \citep{long03} and informal \citep{fine77} settings. Inequalities in cultural access can then be understood as reflecting inequalities in the ability and opportunity to engage in certain forms of informal interaction that Simmel \citeyearpar{simmel49} referred to by the name of ``sociability.''  I argue that the primary uses of aesthetic consumption are social \citep{dimaggio87}.  Even if the ``gratifications'' that the consumption of aesthetic goods provides can sometimes be individual, private and  \citep{denora00, hennion01, holbrook_hirschman82}.

This has important implications for our understanding of the link of culture and stratification in late-modern societies.  If the argument offered here is on the right track it becomes clear that for culturally-advantaged members of high-status groups the deployment of cultural capital {\em in interaction} \citep{eliasoph_lichterman03} during face-to-face talk serves as the primary {\em relational mechanism} \citep{tilly98} that links social position to aesthetic consumption.  This mechanism explains both status-based advantages premised on the maintenance of (homo)social bonds, as argued in the classic Weberian tradition, as well as those social advantages tied to the ability to connect across distant social positions in the social structure \citep{erickson96, lizardo06a}.

\section{Culture as a social resource:  Specifying the phenomenon}
The empirical issue that is the focus of this paper pertains to the ``behavioral implications'' of the aesthetic consumption of popular (and fine arts) culture \citep{fine77}.  This phenomenon can best be characterized through the frequently noted observation that cultural knowledge obtained from aesthetic consumption experiences will become the primary content used for the initiation and maintenance of conversational interaction \citep{dimaggio87, collins75}.  I refer to this (presumably) pervasive phenomenon as {\em culture-mediated social interaction}. Culture-mediated interaction can technically defined as the deployment of cultural knowledge obtained from aesthetic consumption experiences as a resource to generate conversation. In a nutshell (and less technically): the discussion of ``cultural matters'' with friends.

Culture-mediated social interaction is a primary example of ``sociability'' in late-modern societies, and a key marker of social status \citep{collins75, lizardo06a}.  Sociability is social interaction that is not driven by explicit instrumental purposes, but that is keyed to the continuing repetition of episodes of social intercourse for their own sake \citep{simmel49}. This is a form of social intercourse explicitly dissociated from direct instrumental pursuits, as a form of association ``which does not have a strictly economic or business purpose'' but access to which (and restrictions on) are a key marker of the status group boundaries \citep[114]{weber94}. It is in this respect that sociability can be thought of as the ``play form'' of the more general process of association (coming together from some purpose).  This interaction cannot be driven by traditional ``instrumental'' contents (e.g. work, politics, etc.), but requires instead ludic content that serves to simply sustain the focus of attention in the interaction itself \citep{collins75}. I argue that the ``entertainment talk'' typical of those who are active culture consumers is one of the forms of cultural knowledge that is most apt for these playful purposes.
        
This Simmelian mechanism is at the center of Collins' microinteractionist theory of culture and stratification.  According to Collins, 
        
\begin{quotation}
[w]ith conversation for enjoyment, we have left the serious realms and are now fully inside the world of sociability.  The social relationships involved also change.  Practical talk establishes relationships for a very limited purpose \ldots alliances based on extrinsic rewards, not on personal liking \ldots Entertainment talk operates on a different level.  It is talk for the sake of talking, as an enjoyment in itself.  Here the talk is embedded in nonverbal rituals that are even more important than the words themselves \ldots People who participate in it to any substantial degree are friends of some sort \ldots \citep[124]{collins75}.
\end{quotation}

Following this line of reasoning we can begin to appreciate with clarity the ``social'' function of symbolic goods consumption. These functions are intimately related to the maintenance and enhancement of status-linked social cohesion in small groups \citep{fine77}.  In this respect, Fine's basic complaint in relation to research of popular culture in the early 1970s, can be repeated vis a vis research on cultural stratification today.  This research, generally ignores ``the behavioral implications of \ldots culture [consumption], particularly as these implications relate to social interaction'' \citeyearpar[453]{fine77}.  Thus, if culture results in talk, and the micro-realities of stratification and status are reproduced through talk, then culture is intimately implicated in the everyday, mundane reproduction that most researchers interested in stratification dynamics study from more Olympian heights \citep{collins04}.  

It is precisely because ``culture produces considerable discussion, and thus \ldots[its] content \ldots will become a resource in conversational interaction'' \citep[453]{fine77} that understanding the phenomenon of culture-mediated interaction should be a primary goal of those interested in the linkage between culture, stratification and segmentation of social interaction by status.  Accordingly, those cultural goods which find an audience, will be deployed as the basis of conversation and will read the ``fodder for least-common denominator talk.'' We should expect that the consumption of aesthetic goods infuses local conversation-based interaction rituals with a common focus of attention allowing persons to spend time engaged in (Simmelian) interaction for its own sake \citep[443]{dimaggio87}. For instance, Elizabeth Long \citeyearpar{long03} in her study of women's reading groups in Houston, Texas, finds that cultural knowledge and taste become the primary content of social interaction. Consistent with the claim that connects the consumption of aesthetic products and sociability, she finds that reading group members ``tend to press books into service for the meanings that they transmit and the conversations they generate \citeyearpar[73]{long03}.  

It is safe to conclude, in agreement with an earlier tradition on the behavioral consequences of media and the arts, that the consumption of cultural goods has many ``effects.'' But these effects go beyond its perceptual, psychological and emotional impact at the point of direct reception.  The main effect of culture is social, in terms of producing content that becomes ``processed'' in face-to-face interaction.  We can then think of ``[t]he interacting group \ldots as a culture processing machine, which takes popular culture as experienced by group members, and transforms it into content useful for the continuation of interaction'' \citep[460]{fine77}.  I propose that in the contemporary context, aesthetic experiences serve the role of providing fodder for ``entertainment talk'' \citep{collins75, dimaggio87, fiske87, frith98}, and that the primary uses to which individual put the knowledge gained through experiences of aesthetic consumption are  ``social'' in the aforementioned sense.  

Thus characterized, the key empirical dynamic that is the focus of this paper concerns the phenomenon of conversation that involves references and allusions to arts and culture as its primary content.  We should expect that arts consumption will be associated with a higher frequency of these types of culture-mediated social interaction. In addition we should also expect that the relative frequency of these opportunities for culture-mediated interaction will be stratified along the same lines that culture consumption and social connectivity are stratified \citep{dimaggio87, lizardo06a}.  This means that individuals who are already have a higher predisposition to engage in aesthetic consumption---indexed by levels of schooling \citep{bourdieu84, bourdieu_darbel91}---will also be more likely to themselves utilize the cultural knowledge gained from this experience to engage in social interaction.

\section{Theoretical and Substantive Importance of the Phenomenon}       
 Why should the phenomenon of culture mediated interaction be of interest to sociologists interested in the linkage between culture and stratification?  While the connection between aesthetic consumption and social interaction (and thus the uses of culture to build and maintain social networks) has been a relatively neglected issue in the current sociology of culture literature the phenomenon has been obviously perceived to be a key issue in some of the foundational theoretical statements in the field. From this literature we can garner a series of leads as to why we should find that culture-mediated interaction is relatively generalized and why it should be stratified by the familiar markers of status.  
                   
\subsection{Relationship formation}
Culture-mediated interaction allows for the formation and renewal of social relationships.  This is a ``need'' that is structurally induced by the fact that, as recent research on the dynamics of personal networks show, new relationships are constantly being formed and old ones renewed (or dropped) during the adult life-cycles.  The most surprising thing about ``intimacy'' and ``friendship'' is thus, how unstable it tends to be \citep{wellman_etal97}.  Cultural knowledge gained through aesthetic consumption is a key resource with which to form {\em new} social relationships.  ``Culture talk'' is especially useful at the \emph{initial} stages of ``the acquaintance process'' \citep{newcomb61}. Culture-mediated interaction is helpful in this respect because it allows us to search for ``deeper'' areas of commonality and agreement without much risk (as would inquiring about politics or economic position), so that the interaction can proceed smoothly.  Culture mediated interaction is a lubricant that keeps conversation going when intimate content is not yet available.   When individuals do not have a previous social or cultural bases for interaction based on common membership in the same social network or common biographical experiences they instead rely ``\ldots upon common cultural topics on which to build a conversation.''   Cultural goods and aesthetic experiences in general  provide ``\ldots a topic which is both interesting and non-controversial.  It provides an entr�e to interaction'' \citep[460]{fine77}. 

Culture-mediated interaction that makes use of cultural knowledge about symbolic goods that are truly popular may cut across class and status lines may in this way serve to connect persons across different positions in the social structure and across different taste publics, facilitating coordination across hierarchical positions \citep{erickson96}. It is precisely because ``[t]he content and symbolic meaning of these products has diffused sufficiently widely as to make them a viable conversation topic for strangers and new acquaintances'' \citep[460-461]{fine77}, that it  allows persons to create and sustain social connections--even if sometimes brief and sporadic (so-called ``weak ties'')--when there exist large socio-demographic differences between them \citep{lizardo06a}.

\subsection{Status-based solidarity in small groups}
Culture-mediated interaction serves as the primary mechanism to generate and reproduce the micro-solidarities that hold status-groups together \citep{ostrower98}.  The common knowledge gained through culture talk serves as a cognitive and emotional basis of solidarity within the micro-ritual generated in the interaction, as persons become entranced by ``the joy of sharing names'' \citep{dimaggio87}. Culture-talk``\ldots facilitates the sociable intercourse necessary for acquaintanceships to ripen into friendships'' \citep{dimaggio87}.   

Status-based exclusion occurs largely through micro-mechanisms governing entry (or determining systematic exclusion from) class-structured micro-worlds organized around specific conversational interactions \citep{collins75, collins00}.  The experiential and phenomenological ``stuff'' of stratification is essentially based on talk, and much of this talk is about culture.  This is the reason why we can consider taste to be ``\ldots a form of ritual identification and a means of construction social relations\ldots It helps to establish networks of trusting relations that facilitate group mobilization and the attainment of such social rewards as desirable spouses and prestigious jobs''  (DiMaggio 1987:  443).  Culture-mediated interaction thus serves to increase the cohesion and social closure of small groups, providing a way to segment them from the larger social structure;  this is why we can think of culture as being much more ``\ldots than just an acquaintanceship device.  Member know each others' backgrounds, and have a general idea of others positions, values, and beliefs.  In this case popular culture can serve as a mechanism for establishing, confirming and reifying group beliefs and attitudes'' \citeyearpar[461]{fine77} 
                   
The primary function of aesthetic consumption experiences is precisely to serve as the non-instrumental content that justifies the existence of status-segregated conversational rituals:  ``\ldots talk about entertainment-the discussions of movies, music, sports, actors and athletes\ldots make up so much of casual chitchat.'' This becomes even more important ``\ldots if we include talk about things people consume or seek out-food, clothes, houses and gardens, restaurants scenery, vacations.  In some social circles, this entertainment talk makes up most of the content of conversations; in all classes, it is important" \citep[125]{collins75}.  The Simmelian language of ``social circles'' is not accidental, since as Kadushin \citeyearpar[797]{kadushin66} noted in his classic study on the subject, Simmelian urban social circles are produced and reproduced through shared {\em cultural tastes}--e.g. concert going, visiting the museum, etc.--and other activities intrinsically connected to culture-mediated interaction.                     
                   
\subsection{Linkages within loosely bounded and sparsely knit networks}     
Culture-mediated interaction is not only  the ``bond'' that keeps status groups together, but also the ``solvent'' that allows for sporadic connections across status lines and that provides new forms of high-status group solidarity not built around traditional notions of commonality anchored to geography, community and shared institutional affiliations and kinship structures \citep{holt98}.  One of the most significant transformations in the social logic that governs status in contemporary (post)industrial societies consists of the rising role of the arts, and mass produced culture in general, in providing the ``baseline'' forms of cultural knowledge necessary to sustain interaction for persons who cannot come to expect commonalities based on traditional symbols of acquaintance and recognition. 

This shift has occurred in the context of the social transformation of status-based collectivities from solidary groups with a strong allegiance to community and place, to looser, wide-spanning ``personal networks'' characterized by increasing geographic and social mobility \citep{wellman_etal88}.  This process acquires more importance as these network ties have been transformed in the contemporary context of increasing geographic mobility \citep{griswold_wright04} and the decline of the traditional bonds characteristic of traditional local communities \citep{wellman79}, which used to be centered around kinship and spatial contiguity \citep{dimaggio_mohr85}.  As the personal networks of socially mobile members of the upper-middle class have become more fluid, loosely knit, more discontinuous and less tied to spatial contiguity and kinship \citep{wellman_wortley90, wellman_etal88}, the arts and popular culture have come to increasingly serve as the ``default" forms of knowledge that connect people across different interaction foci.
                   
Social interaction based on culture comes to occupy center stage because aesthetic consumption experiences are an ideal ``content'' to fuel the process of Simmelian sociability. In contrast to material goods which ``are physically present and visible'' and which are more likely to produced invidious distinction, cultural consumption, ``\ldots is invisible once it has occurred.  This evanescent quality makes artistic experience, described and exploited in conversation, a portable and thus potent medium of interactional exchange.'' \citep[443-444]{dimaggio87}.  Popular culture and the arts (along with ``fashion, cuisine and sport'' are the only analog to a common cultural currency in complex societies \citep[443]{dimaggio87}.  The arts are not only a generator of cultural knowledge and are also one of the most important facilitators of culture-mediated interaction.                 
            
In sum, some of the most important lines of theory in the culture and stratification literature point to the central importance of the phenomenon that we have labeled and characterize as {\em culture mediated social interaction}.  These theoretical streams point to this phenomenon as a pervasive, ubiquitous feature of late modern societies.  According to this line of theory, culture-mediated interaction is key to understanding the dynamic, interactional origins of the social discontinuities and status-based social boundaries that characterize the modern social order \citep{illouz97}. This makes its almost total neglect in the more empirically oriented recent literature in culture and stratification an even more puzzling oversight. 
                   
The formal and informal system of production of aesthetic goods provide as ``input'' the very content that serves as the scaffolding around which everyday episodes of ``talk'' and sociability are built.  The socially segmented micro-orders based on face-to-face (and increasingly electronically mediated) interaction require a constant stream of new interactional content in order to be sustained and thus reproduce themselves over time. It is thus no surprise, as we saw above, that the stratification of culture-mediated interaction should mirror in non-trivial respects patterns of stratification found in the consumption of aesthetic goods and experiences. This is consistent with Simon Frith's \citeyearpar[4]{frith98} suggestion that ``[p]art of the pleasure of popular culture is talking about; part of its meaning is talk �''   The micro-solidarities that form the basis of status-linked orders but also form of the core of what has come to be called ``social capital'' are may thus be built around the most seemingly innocuous forms of conversational subject matter.  This implies that the phenomenon of culture-mediated interaction is relevant for those interested in the connection between cultural and social capital \citep{lizardo06a, jeannotte03, dimaggio96}.  As John Fiske's observed, civic membership and ``belonging'' to neighborhoods and communities is based on talk, and ``[m]uch of this talk is about the mass media and its cultural commodities'' \citeyearpar[73]{fiske87}.
                
\section{Empirical Analysis}
\subsection{Data Source}
In what follows I examine data taken from the Arts and Culture Supplement to the National Social Survey fielded in 2002 and 2004.   The NSS is a weighted (by age, education, marital status, gender, race, ethnicity) probability sample respondents obtained through computer-aided telephone interviews. This survey is particularly appropriate for our purposes because it contains information not only about culture consumption experiences, but also about the likelihood that individuals engage in informal social interaction featuring cultural content taken from those aesthetic consumption episodes.  Because the items of interest were fielded in both waves, I pool the samples for which I have complete data on all socio-demographic variables from the 2002 ($N=955$) and 2004 ($N=776$), which leaves us with a total combined sample of 1,731 respondents. 

\subsection{Culture-Mediated Social Interaction Survey Items}

As part of the survey, respondents were asked the following question:

During the past twelve months, approximately how often have you had a discussion with someone you know about any of the following:
\begin{enumerate}
\item  {\em A book, a poem, or a story?}
\item  {\em A musical performance, live or recorded.}
\item  {\em A movie?}
\item  {\em A play or other dramatic performance?}
\item  {\em An art show or work of art?}
\end{enumerate}

Respondents were asked to choose from the following response categories: 1) very often, 2) fairly often, 3) not very often, 4) or not at all. 
                
\subsection{Aesthetic Consumption Survey Items}
                                            
Respondents were also asked to report the frequency with which they engage in various forms of aesthetic consumption:
 
\begin{enumerate}
\item  {\em Books}:  Approximately how often during the past 12 months, have you read a book?
\item  {\em Movies}:  Approximately how often during the past 12 months, have you been to the movies? 
\item   {\em Music}: Approximately how often during the past twelve months would you say you listened to music at home, for instance on CDs, tapes, records, a computer, or on the radio?
\item  {\em Live performances}: Approximately how often did you go to any live performances such as plays, concerts, musicals, or dance?
\item  {\em Art}:  Approximately how often did you go to any art shows or museums during the past 12 months? 
\end{enumerate}  

\section{Analysis and Results}
\subsection{Distribution of culture-mediated interaction in the United States}
\subsubsection{Overall frequency}
I begin by attempting to provide an answer to a simple question: what is the relative distribution of the frequency of culture-mediated interaction in the contemporary United States? Insofar as culture-mediated interaction is a key process in current theories of the culture/stratification link, we should expect it to a fairly widespread phenomenon.  Because there has been virtually no quantitative research on issue however, it is hard to estimate with certainty how much is "widespread."  Given what we know from the few qualitative studies that have been conducted, we can  expect to observe close to a majority of the population reporting high or moderately high frequencies of culture-mediated interaction.\footnote{As with any study that uses arts participation data in order to draw larger conclusions about the operation of culture in society at large, the big caveat is of course that any observed frequency of culture-mediate interaction is probably an underestimate, as persons can only report interacting with friends on the basis of the cultural forms included in the survey.  Clearly, persons who are observed to not interact with friends very often on the basis of content taken from the cultural domains included in the survey, may still be doing so on the basis of cultural forms not included in the survey.}
      
Figure~\ref{fig1a} shows the univariate distribution of culture-mediated interaction across the five broad realms of cultural activity included in the 2002 and 2004 National Social Surveys.\footnote{the circle is the observed percentage and the vertical line represents the 95\% confidence interval, calculated using the ``exact'' binomial estimate for the observed proportion.}  The distribution of culture-mediated interaction frequencies across all five cultural domains supports the expectation that persons use culture as a discussion topic in a fairly regular manner.  78\% of respondents report frequently (either ``fairly'' or ``very'' often) using {\em at least one} of the cultural domains included in the survey as a subject matter in conversation with friends.
                
When it comes to the three most popular cultural domains  (movies, literature and music), it is clear that a majority (or close to a majority) of respondents report having discussed matters related to that cultural domain with a friend.  About 60\% of persons report talking to friends about movies ``very often'' or ``fairly often''; the corresponding percentages for books and music are 52\% and 46\%. Even for the two cultural domains which are expected to have the most exclusive audiences (in terms of such socio-demographic characteristics as education and age), in this case theater and the arts, 26\% and 19\% of respondents report relative frequent rates of culture-mediated interaction respectively. 
                
\subsubsection{Frequency among cultural active respondents}
The raw estimates of culture-mediated interaction give a somewhat misleading picture of the extent to which the phenomenon occurs. It is reasonable to suppose that a necessary precondition for persons to deploy any type of cultural knowledge during interaction is that they have acquired this knowledge.  It is a good bet that knowledge about symbolic goods is acquired through direct engagement with the formal and informal culture industries (and non-profit organizational fields) that deliver these goods.  Thus, we should expect the relative chances of observing high rates of culture-mediated interaction to be mediated by the extent to which the persons engages the cultural field in question as a consumer.
      
Figure~\ref{fig1b} shows the conditional univariate distribution of culture-mediated interaction among those respondents who report engaging in the corresponding cultural activity either ``fairly often'' or ``very often.''  I find that not surprisingly, rates of culture-mediated interaction among active consumers of the cultural forms in question are substantially higher than the average.  About 90\% movie, 71\% of literature and 51\% of music consumers report engaging in relative frequent rates of culture mediated interaction.  The most dramatic relative increases in relation to the average is observed for the two cultural domains associated with the non-profit ``fine arts'' field:  53\% of active dramatic theater consumers and 61\% of active art consumers report discussing this type of cultural content with a friend.  
      
Interestingly, the univariate distribution for culture-mediated interaction based on music shifts the least after consumption is taken into account.  This indicates that this type of cultural content is deployed routinely without necessarily being tightly connected to consumption.  It can also mean that the ``uses'' of music as cultural form may be unique:  music may constitute both a ``social'' resource for the cementing of connectivity in social relationships, and also a private ``resource'' for personal enjoyment and the customization of ``life-spaces'' \citep{denora00}.
      
Another noticeable fact is that both movies and literature seem to be ``inherently social'' cultural domains in these data.  It thus appears to be almost (statistically at least) ``non-sensical'' to be an active movie consumer or book reader without using this cultural knowledge in interaction.  The concrete mechanisms governing the way in which these two cultural forms are deployed in everyday settings may be distinct however. The fact that in these data the proportion of respondents who actively engage the cinematic field {\em without} at least occasionally using this cultural knowledge in conversation corresponds well with Fine's \citeyearpar{fine77} original proposal that cultural goods with widespread appeal would also be the most likely to be mobilized in interaction in small groups. Research on book groups by Long \citeyearpar{long03} and Illouz's \citeyearpar{illouz03} research on Oprah's televisual ``book club'' suggest that books are a primary conversational resource among a wide variety of audiences.  However, literature conversations tend to occur within the relatively formalized ambit of the formal groups purposefully designed for the purposes of literary discussion (for a history of the book club institution in the U.S. see Long \citeyear{long03}).  This is in keeping with the highly institutionalized perception of (most) literature (and reading as a cultural practice) as ``serious.''   Conversations about such cultural fare as movies, however, serves to cement interaction in less formalized social settings \citep{fine77}.
      
\section{The stratification of culture-mediated social interaction}
While most persons are indeed likely to engage in some form culture-mediated interaction, it is to be expected that as a general rule, culture-mediated interaction should be stratified along indicators of status and prestige.  We should expect this to be the case for two reasons:  first, culture consumption is itself stratified, with more educated individuals being more likely to engage in a wide range of aesthetic consumption activities than the less educated \citep{bourdieu84, bourdieu_darbel91, vaneijck99}, and as we have already seen, culture-mediated interaction is much more likely to be observed among active culture consumers than among inactive ones.  
                
Second, opportunities for social interaction and such indicators of social connectivity as network size and proportion non-kin in the typical discussion networks are stratified by the same types of status and prestige indicators, in particular educational attainment \citep{marsden87, fischer82, burt90}.  Insofar as culture mediated interaction is more likely to take place among members of informal cliques, we should expect that those persons most likely to have wide-spanning, loosely knit networks of acquaintance, composed of elective relationships of ``friendship'' (in contrast to ascriptive relations of kinship) should be more likely to engage in culture-mediated interaction.  Most research on personal discussion networks indicates that the higher educated are more likely to enjoy these types of local relational environments than the less educated \citep{burt90, marsden87, erickson96, lizardo06a, fischer82}.      

\subsection{Estimation}
The key challenge in estimating the effects of socio-demographic markers on the propensity to engage in culture-mediated interaction is tied to the peculiar nature of this dependent variable.  More specifically, the effect of any one socio-demographic characteristic on culture-mediated interaction in any one domain is likely to be shaped by the fact that the distribution of individuals along a culture \emph{talk} continuum is simultaneously determined (in ways that we will attempt to disentangle below) by their distribution along a culture \emph{consumption} gradient (and vice-versa). Because it is likely that the same predictors of culture talk will also figure as predictors of culture consumption, the issue of the mutual determination of talk and consumption has to be dealt in any attempt to model the socio-demographic heterogeneity in the propensity to engage in both activities.  

\subsection{Statistical model}
I specify the effects of social position on the propensity to engage in culture-mediated interaction using a binary ordered probit model \citep{ greene98}. This model is useful precisely when the analyst if faced with limited dependent variables which may share a causal history or are simultaneously determined \citep{calhoun89}.  In the standard univariate probit specification for instance, the underlying propensity for each person to engage in culture-talk or aesthetic consumption would be conceptualized as an unobserved latent (normally distributed) variable ($y^{*}$).  The empirical realization ($y$) of the outcome variable is binary and equals one if the underlying latent propensity is above some threshold (e.g. $y^{*} \geq 0$) and zero otherwise. The probability of being above or below the threshold is stipulated to be a linear function of a vector of socio-demographic covariates ($\beta x$), and all the other factors that affect the outcome not contained in this vector, are assumed to be independent and identically distributed.   

The general form of the bivariate probit model is shown in Figure~\ref{fig3a}. The bivariate probit takes the same general form as the univariate probit model just described \citep{long_freese06} (models (1) and (2) in Figure~\ref{fig3b}), with the exception that the two dependent variables are now considered to be the empirical realization of two latent underlying variables ($y^{*}_{1}$ and $y^{*}_{2}$) presumed to follow a bivariate normal distribution.  This model is usually referred to as the ``seemingly unrelated'' bivariate probit. The relationship between the latent normal variates ($y^{*}$) and the observed binary variables ($y$) is as follows:
\begin{equation}
y=g(y^{*})=\left\{\begin{array}{rl}
                                0 &\mbox{if $y^{*} \leq 0$} \\
                                1 &\mbox{if $y^{*} > 0$} 
                        \end{array} \right.
\end{equation}
The seemingly unrelated bivariate probit model is then given by:
\begin{equation}
\begin{array}{rl}
y^{*}_{1i}=\beta_{1} x_{1i}+u_{i} \\
y^{*}_{2i}=\beta_{2} x_{2i}+v_{i} \\ \\
corr(u,v)=\rho \\
u,v \approx \Phi(0,1) 
\end{array}
\end{equation}
Where $\Phi$ is the cumulative normal distribution function, with mean zero and standard deviation of one. The model allows for the correlation ($\rho_{12}$) between the exogenous errors corresponding to each probit equation ($u$ and $v$) to be estimated from the data. 

In the current application, this parameter represents the \emph{shared} exogenous factors that affect both culture talk and culture consumption and which are not included in the systematic component of the model. In this way, the coefficient estimates corresponding to the predictors of culture talk can be estimated net of the association of culture-mediated interaction with those factors that also affect culture consumption.

\subsection{Variation in culture-mediated interaction by social position}
The coefficient estimates (and the associated t-statistic corresponding to the hypothesis that the regression weight corresponding to that variable is zero) for each social position indicator are shown in the left-most set of columns (for each cultural domain) in Tables~\ref{tab1} and ~\ref{tab2}.  A t-value with an absolute value above 1.96, indicates that members of that category have a statistically higher or lower (depending on the sign of the coefficient) underlying propensity (in comparison to  average respondent) to engage in culture-mediate interaction. Likelihood ratio $\chi^{2}$ statistics corresponding to the test of the hypothesis that a given block of variables (e.g. education, age) is equal to zero, are shown in the top panel of Table~\ref{tab3}. The last three columns show the proportion of the total model $\chi^{2}$ that each separate block of variables is responsible for.  These percentages give a rough indication of the relative importance of that set of variables (in comparison to one another) in predicting the propensity of culture-mediated interaction. 

The bivariate ordered probit equations shown in Tables~\ref{tab1} and ~\ref{tab2} have all-categorical predictors in the right-side.  To make interpretation of the relevant effects more manageable, I use \emph{contrast} coding instead of the usual dummy-coding for all predictor variables.  This means that all of the categories (for those variables that have more than two categories) are included in the model. In these coding scheme, there is no omitted reference category since the coefficients are constrained to sum to zero across all of the different values of a given variable. For variables with only two categories only one coefficient is shown, since the other one has the same absolute value but opposite sign. In contrast coding, a statistically significant coefficient estimate corresponding to a given category (e.g. high school education) has a simple substantive interpretation:  it indicates that members of that category are more or less likely to have a larger propensity to engage in culture-mediated interaction in that domain \emph{in comparison to the sample mean across all of the categories} (dummy-coding in contrast compares a given category to some arbitrarily omitted one, and thus is more difficult to assign a substantive meaning to the coefficient estimates).   

As shown by the magnitude and direction of the coefficient estimates in Tables~\ref{tab1} and ~\ref{tab2} and the likelihood-ratio statistics in Table~\ref{tab3}, across all five domains, culture-mediated interaction is unequally distributed by education with the higher educated being more likely to use culture in conversation than the less educated ($p<0.01$) \citep{collins75}. Naturally, education does not equally affect culture-talk in all domains.  Differentiation by education (in comparison to heterogeneity premised on age, gender and race) is preponderant in the case of culture-mediated interaction involving music and literature and the arts, and is relatively less important in the case of movies ($\chi^{2}=43.3, p<0.01, df=4$). In the case of the arts, education accounts for about 66\% of the total model $\chi^{2}$, while within the domain of literature education is responsible for about 52\% of the total model fit.

Differentiation according to age categories in the underlying propensity to use culture in interaction, where it exists, favors the young over the old, with younger persons being more likely to mobilize culture as a resource in conversation than older persons.  Age heterogeneity  in the propensity to engage in culture-mediated interaction is very salient for music ($\chi^{2}=61.6 p<0.01, df=6$), but it is \emph{extremely} salient in the case of movies ($\chi^{2}=203.9$), accounting for about 70\% of the total model $\chi^{2}$ in the culture-mediated interaction equation.  This pattern of results makes sense in the context of previous research that shows that the young are more likely to be involved in the cultural domains most likely to be associated with the commercialized arts in comparisons to those embedded in the non-profit arts sector.  Differentiation by age in the propensity to use culture in interaction is only marginally discernible in the plastic and arts ($\chi^{2}=11.5, p=0.08$, and is statiscally significant but relatively weak in literature ($\chi^{2}=14.6$).  Age effects however are somewhat discerinible for the performing arts ($\chi^{2}=14.6, p<0.01$), where members of younger age categories show a higher propensity to engage in culture-talk than older respondents.   

Gender differentiation in the underlying propensity to use culture in interaction (favoring women over men) can be observed in all domains except for music ($\chi^{2}=3.0, p=0.09, df=1$), but the gender difference is relatively weak when it comes to the arts ($\chi^{2}=16.3$) and is most salient in the domain of pleasure reading ($\chi^{2}=66.5$) \citep{tepper00, long03}, with gender accounting for about 39\% of the overall $\chi^{2}$ corresponding to the three variables considered here.  

Racial and ethnic differentiation in culture-talk propensities are generally weak across all aesthetic domains, with self classified blacks and hispanics being about equally likely to engage in culture-mediated interaction involving the arts once we adjust for education.  The only discernible difference appears in the case of the performing arts ($\chi^{2}=7.1, p=0.03$), where I find that persons who self-classify as African American or black are \emph{more} likely to use cultural resources from this aesthetic domain in interaction.  This result may appear counter-intuitive until we recall that both Chaves \citeyearpar{chaves04} and Wuthnow \citeyearpar{wuthnow06} show that one of the primary roles of religious organizations in the everyday life of congregants is as sites in which persons are exposed to the arts, with a disproportionate amount of congregational time and activities devoted to what would be considered the consumption of (relatively high-status)  symbolic goods delivered by non-profit arts organizations. Insofar as congregations are more likely to be a central part of the lives of African American, it is not that surprising that the aesthetic experiences to which they are exposed in those social foci \citep{feld81} come to structure their everyday interaction in informal contexts.  


\subsection{Variation in aesthetic consumption by social position}
How does variation across levels of education, age and gender on the propensity to engage in culture-mediated interaction compare to heterogeneity across the same three social position indicators on the frequency of aesthetic consumption?  For each cultural domain, the second set of columns of Tables~\ref{tab1} and ~\ref{tab2} shows the coefficient estimates for the component of the bivariate probit model predicting the frequency of aesthetic consumption, using the same socio-demographic covariates used to predict culture-based talk. The results show that aesthetic consumption is largely structured in similar ways in comparison to culture-mediated interaction (with important differences to be addressed more fully below), with consumption heterogeneity across education, age and gender, and race being the most important predictors of underlying propensities to engage literature, popular culture and the arts.  The bottom panel of Table~\ref{tab3} shows estimates the relative importance of these four sets of social position markers for participation in the five cultural domains already considered.  

Consistent with almost all studies in cultural participation research, educational differentiation in the underlying propensity to appropriate symbolic goods exists in all domains connected to the non-profit arts, and favors the educated (as frequent over consumers) over the less educated (as abstainers). When it comes to relative rates of engagement engagement in the plastic arts (73\%), the performing arts (56\%) and pleasure reading (56\%), education accounts for the bulk of the total explanatory power in the model.  This suggests that education is the primary source of social differention in the more ``artistically'' oriented domains, but not in the more ``commercial'' domains \citep{bourdieu84}.  In the case of frequency of music listening ($\chi^{2}=1.3, p=0.87$) heterogeneity by educational grouping is statistically non-existent in these data.  For movies, there exist some non-negligible differentation by education ($\chi^{2}=46.6, p<0.01$) however. 

The effects of age in terms of differentiating audiences across domains largely mirror its effects in differentiating persons according to their underlying propensity to use culture in interaction. In the domains of music and movies age apears as the primary factor differentiating avid consumers from abstainers ($\chi^{2}=84.3, df=4$ for music and $\chi^{2}=176.0$ for movies), with the young being more likely to engage the popular arts outside the home than the old.  There is no comparably strong differentiation across age-groups in terms of frequent engagement in the domains associated with literature and the non-profit arts, once we hold constant education, although a relatively weak age-differentiation effect can be ascertained for the performing arts ($\chi^{2}=19.7, p<0.05$) and for literature ($\chi^{2}=13.0, p<0.05$).  

The effect of gender as a differentiating factor in arts participation frequency also mirror the gender effects on culture-mediated interaction, with the strongest gender effect being restricted to reading ($\chi^{2}=33.6, p<0.01$) \citep{tepper00, long03}.  Differentiation by gender in consumption propensities is largely negligible in the plastic and performing arts ($p>0.25$), although there is a slight ($\chi^{2}=7.8, p=0.02$) gender effect in the case of music (with women being relatively more likely to be frequent music listeners than men), and even more negligible effect in the case of movies ($p=0.09$).   

\subsection{Variation in the effects of social position across the consumption/interaction divide} 
As should already be clear, while social differentiation in the propensity to engage in culture-mediated interaction largely mirrors that found in terms of the likelihood to actually be an active consumer in those domains, there are also substantively relevant discrepancies.  Table~\ref{tab4} shows the $\chi^{2}$ statistics of a series of likelihood ratio tests comparing the bivariate probit models shown in Tables~\ref{3}-~\ref{4}, with constrained versions of themselves in which the coefficients for Education, Age and Gender and Race are fixed to be the same in both the culture-talk and culture-consumption equations.  Essentially this is a statistical test that the effects of these variables are the same for both culture-mediated talk and frequency of consumption in each domain.  A high value of the $\chi^{2}$ statistic and a \emph{p}-value below conventional levels indicates a statistically significant difference between the unconstrained and constrained models, suggesting that we should reject the null hypothesis of equal effects of that social position marker for culture-talk and culture-consumption.

The table shows that education has similar effects on talk and consumption on all domains except for music, for which the model fit significantly deteriorates when the education effect on culture-talk is constrained to be equal to its effect on consumption frequency ($\chi^{2}=24.3, p<0.01$).   The reason for this, as we have already noted, is because there is no statistically discernible heterogeneity across educational categories in frequency of music \emph{consumption}  (thus music mirrors the more commercial realm of movies in this respect), even though the highly educated are more likely to engage in culture-mediated \emph{interaction} using musical culture as a resource (which means that music mirrors the higher-status non-profit arts and reading in this respect). I also find a weaker deterioration in model fit in the case of arts consumption ($\chi^{2}=9.45$), suggesting that the effects of education in segmenting active art audiences might indeed be relatively stronger than its effects in segregating those groups who use the arts as a resource in interaction ($p=0.05$).

How does age differentiation in terms of culture-talk propensity compares to age-differentiation in cultural participation? I find that age has comparable differentiating effects across the consumption/interaction divide for the domains of Music, Arts and Drama ($p>0.10$). But there exist heterogeneous age effects in the reading  ($\chi^{2}=36.6, p<0.01$) and film  ($\chi^{2}=12.8, p<0.05$) domains. As shown in Table~\ref{tab2}, the key age-differentiation discrepancy in terms of engagement in the culture of the reading class obtains for the oldest age group (65+), members of which are statistically \emph{less} likely than average to engage in culture-talk based on books ($\beta=-0.18, t=-2.98$), but are discernably \emph{more} likely than the average respondent to be frequent readers ($\beta=-0.19, t=3.44$). This suggests that a disproportionate share of that (relatively small) percentage of respondents who report reading regularly but not talking regularly about books (see Figure~\ref{fig1}), may in fact be older persons who maintain high-consumption habits in the literary domain but whose (informal, non-kin) networks begin to experience shrinkage and decay. In the case of music, we find the opposite age contrast as that which obtains for reading: the two oldest age groups (fifty or older) are much less likely to listen to music than they are two talk about it, with generate an asymmetry in the age effect across the two outcomes.  

In all domains except for music and the reading ($p>0.05$) the strength of the differentiation by gender in culture-mediated interaction is more salient than that which pertains to culture consumption.  This means that for movies and the arts gender is a more useful marker in determining who uses culture in interaction than in determing who engages in cultural consumption.  This accords with recent research that suggests that common interests in cultural matters constitutes a key factor facilitating solidarity within gendered workplace (and other informal) networks \citep{lizardo06b}.  

Finally, the differentiating effects of race and ethnic minority classification are comparable in the case of music and the plastic arts.  The model that constrains the race and ethnic effects to be the same exhibits a marginally worsening of fit in the case of literature ($\chi^{2}=5.15, p=0.08$) and the performing arts ($\chi^{2}=5.09, p=0.08$). The most statistically salient disparity in terms of relative differences across racial and ethnic categories in the consumption versus the culture-talk equations  occurs in the realm of movies ($\chi^{2}=8.5, p<0.05$). Table~\ref{tab1} why this is the case:  while Hispanics are no more likely than the average respondent to use film culture as a resource in interaction ($t=-0.56$), they are statistically \emph{more} likely than the sample average to be frequent movie goers ($t=2.84$).  

\section{The relationship between culture consumption and culture-talk:  Two models}
As should already be evident, a key issue that emerges in analyzing the relationship between aesthetic consumption, social position and culture-mediated interaction is the question of their inter-relationship.  This issue is important, because different theoretical representations of the underlying process that generates the association between position, disposition and propensity to engage in culture-mediated interaction imply analytically distinct models of how cultural tastes are shaped by (or shapes) interaction opportunities \citep{lizardo06a, erickson96, mark98b}.  Contemporary research and theory suggests two (parsimonious) ways of thinking about how they are connected to one another.  Essentially this boils down to the question of establishing which of factor is \emph{temporally and causally prior} to the other ones, and which ones are more likely to function as \emph{intervening mechanisms}. 

While it is a commonplace to acknowledge that causal priority cannot be unambiguously established with cross-sectional data, it is mistaken to conclude that the analyst is therefore helpless. Instead, because all hypotheses that suggest some intuitive arrangement of multiple factors put explicit constraints on a given system of inter-relationships (and the relative likelihood that the data were generated by a process consistent with these constraints can established), it is possible to at least ascertain whether the cross-sectional data are compatible to what would be expected by a particular model of the process at ``equilibrium.'' In particular, it is possible to adapt Kendall and Lazarsfeld's \citeyearpar{kendall_lazarsfeld50} influential analysis of spurious correlation for the goal of determining which model is most consistent with a given underlying causal structure \citep{simon54}. In next two sections I briefly summarize these competing views, and propose a set of relatively simple multiple-equation models appropriate for correlated variables with ordered categories that properly characterize them.  These models are shown in ~\ref{fig3b}.

\subsection{Culture-mediated interaction as an intervening mechanism}
One influential substantive model of the process through which culture-mediated interaction is connected to aesthetic consumption conceptualizes \emph{interaction opportunities} in certain delimited social foci \citep{feld81} as the primary (intervening) mechanism connecting social position and aesthetic consumption.  This model is more consistent with a network-theoretic view which sees the residual association between social position and cultural tastes---$\rho_{12}$ Figure~\ref{fig3b}---as a (partially) spurious by-product of differentially distributed (across categorically defined social positions) opportunities for interaction in certain settings. In representing this model (and the taste as intervening mechanism model discussed below) as model (3) in  figure~\ref{fig4} I follow Lever's \citeyearpar{lever79} formalization of Blau's  \citeyearpar{blau55} strategy for determining the most likely cross-sectional association structure consistent with the intervening role of a given variable, but generalize the analytic strategy to the case of two correlated binary variables from within the framework of recursive bivariate probit models.  

Substantively, model (3) in figure~\ref{fig4} says that the observed residual association between social position (and aesthetic consumption should be greatly reduced within levels of social interaction. This model implies that $\delta_12>0$ and that therefore increasing levels of interaction involving culture drive the frequency of culture consumption in that domain. This also means that social position has both direct effects (through $\beta_{1}$) and indirect effects (mediated by $\delta_{12}$) on aesthetic consumption. Thus, those persons who, due to the possession of more expansive and diverse ``personal communities'' embedded in networks or access to social foci in which culture-talk is more likely to occur, are expected to be more likely to acquire and sustain taste for the arts and are also the persons most likely to possess certain markers of status and position in the social structure.

\subsection{Aesthetic consumption as an intervening mechanism}
A second way to conceptualize the generative process that links culture-mediated interaction and aesthetic consumption is to construe \emph{cultural tastes} (as expressed in aesthetic consumption frequencies) as the primary intervening mechanism joining social position and culture-mediated interaction \citep{lizardo06a}. The graphical representation of the model is shown as model (4) in figure~\ref{fig4}.  This conceptualization differs from the network-theoretic model suggested by (3), because here tastes are seen as a much more durable part of the person---``embodied'' cultural capital in Bourdieu's \citeyearpar{bourdieu86} terms---and it is through the differential distribution (across social positions) of (expressed) dispositions towards culture and the arts that we can account for the association between social position and opportunities to engage in culture-mediated interaction. 

Substantively, this model says that a substantial part of the variation in culture-mediated interaction across social positions in the propensity to engage in culture-mediated interaction is explainable by the direct effect of aesthetic consumption ($\delta{21}$ in model (4)) on this outcome.  This model is  consistent with a ``culture conversion'' \citep{lizardo06a} argument regarding the process that links social position with culture-mediated interaction opportunities.  Under this model, cultural tastes serve as ``bias'' in contact selection and contact-maintenance \citep{dimaggio93b, vaisey_lizardo}, such that persons with a given set of tastes tend come to occupy certain social positions (e.g. the rank of the highly educated) and come to surround themselves with contacts of similar-tastes, which then generates opportunities for culture-based talk based on those aesthetic experiences \citep{frith98}. Cultural tastes are thus the intervening mechanism connecting a given social position with a higher likelihood of engaging in culture-talk.

\subsection{Estimation}
In econometric jargon, the model in Figure~\ref{fig3} is the seemingly unrelated bivariate probit.  Models (3) and (4) in Figure~\ref{fig4} are consistent and efficient variants of the SUR model known as the recursive simultaneous equation bivariate ordered probit models.  

The recursive bivariate probit with an endogenous regressor can be written as:

\begin{equation}
\begin{array}{rl}
y^{*}_{1i}=\delta_{21}y_{2i}+\beta_{1} x_{1i}+u_{i} \\
y^{*}_{2i}=\beta_{2} x_{2i}+v_{i}
\end{array}
\end{equation}

These models are especially useful when dealing with the problem of estimating directional effects with simultaneously determined endogenous outcomes \citep{greene98}.  What most analysts usually miss is that in these models there is no need to include additional ``instrumental variables'' or to rely on other so-called ``exclusion restrictions'' to estimate the endogenous effect (this means that the $x_{1}$ and $x_{2}$ vectors can contain the same predictors) \citep{monfardini_radice08}.  Instead, one can simply ``plug in'' one of the jointly determined variables as a right-hand side predictor of the other---as depicted graphically in (2) and (3) \citep{greene98}. The analyst can then use the usual indices of model fit to determine which model gives a better representation of the data-generating process. Including exogenous predictors in $x_{1}$ that are excluded from $x_{2}$, while not strictly necessary for identification, may help make inference more robust to distributional misspecification \citep{monfardini_radice08}.

\subsection{Results}
Table~\ref{tab2} shows the fit statistics corresponding to the models shown in Figures~\ref{fig3} and ~\ref{fig4} for the five cultural domains included in the NSS.\footnote{The consumption component of the recursive bivariate probit model includes additional control variables obtained from a module on ``barriers to arts participation'' added to both the 2002 and 2004 waves.  These questions contain items dealing with possible reasons for why the respondent may be impeded from participating in the arts more often (poor health, lack of time, lack of friends with an interest in the arts). Three additional items deal with whether the respondent's current attendance rates would increase if certain conditions were met (tickets would be cheaper, there was more information available or art venues were closer).  These three items were asked of all respondents.  In addition identical items were asked in the special case of movies and dramatic performances (in the equations predicting these outcomes both sets of variables are included in the analysis). Similar items were asked of music performances and reading in the 2002 but not in the 2004 wave.  Because of this I do not include these items in the consumption equations for these domains.}  For comparison purposes, I also estimate ``naive'' (i.e. single-equation, without correlated errors) versions of the recursive bivariate probit models (3) and (4) in Figure~\ref{fig3}, shown there as models (1) and (2).  These models allow us to ascertain how much of the raw causal effect is a spurious by-product of common unmeasured factors that affect both consumption and interaction. Thus, the key comparison that emerges is between the $\alpha$ and $\delta$ estimates of the effects of culture-mediated interaction on consumption (and vice versa). 

The estimates of $\rho_{12}$ obtained from fitting the bivariate probit models without (baseline) and with (SUR) covariates are shown in the first four columns of Table~\ref{tab3}.\footnote{The seemingly unrelated regression model between two binary endogenous variables without covariates yields the zero-order tetrachoric correlation coefficient between the two interaction and consumption.}  The change in $\rho_{12}$ from the baseline condition to the revised estimate in the SUR specification gives us a sense of how much of the association between consumption and culture-mediated interaction is accounted for by the inclusion of the socio-demographic predictors. For all domains  $\rho_{12}$ declines once we hold constant individual heterogeneity in terms of socio-demographic markers.  However, while for some domains the decline in $\rho_{12}$ is substantial (Music, Books) in other domains it is negligible (the Arts).  

An important thing to notice by looking at the magnitude of the $\rho_{12}$ estimate and the associated p-values corresponding to the test that the $\rho$ parameter equals zero is that the null model of spurious association by social position is rejected for all cultural domains. This means that (as expected) interaction and consumption are not exogenous (i.e. independent) in relation to one another, but may be mutually determined, since there remains a residual association between aesthetic consumption and culture-mediated interaction that is not accounted by socio-demographic heterogeneity.  Thus, it is plausible to consider models (such as models (3) and (4) in Figure~\ref{fig4}) that conceptualize either culture-talk as a mediator---e.g. the ``intervening variable'' \citep{lever79}---between position and cultural dispositions or models in which cultural dispositions are converted (through their linkage to social positions) into opportunities to engage in culture-mediated talk.

How do these two alternative models stack up against one another (and in comparison to the seemingly unrelated specification)?  The fit statistics are shown in Table~\ref{tab2}.  I use the BIC statistic as a measure of model fit. The BIC statistic is an information-theoretic measure that rewards parsimony in model specification \citep{raftery95} and which can be used to select between non-nested models.  The results are easy to summarize:  for all five cultural domains, the aesthetic consumption as intervening variable model (mode (4) in Figure~\ref{fig4}) outperforms the seemingly unrelated and the interaction as intervening mechanisms specifications. In fact, in all five domains, specifying a direct effect going from interaction to aesthetic consumption worsens the model fit over the seemingly unrelated specification (even though this model only uses an additional degree of freedom to fit the $\delta_{12}$ parameter).  This is strong evidence supporting the contention that the data are likely to have been generated by a process consistent with a ``culture conversion'' account (in which pre-existing propensities to consume culture drive the creation of opportunities to engage in culture-talk) than a traditional network model in aesthetic consumption is driven by propensities to interact.

This conclusion is supported by the results shown in the lower panel of Table~\ref{tab3}.  The key comparison there is between the  naive estimate of the effect of the culture-mediated interaction ($\alpha_{12}$) obtained from a simple univariate probit model with aesthetic consumption as the dependent variable, and the corresponding effect in the recursive bivariate probit specification ($\delta_{12}$).  For the estimate of the effect of culture-mediated interaction on aesthetic consumption corresponding to model (1) in figure~\ref{fig3}, we can see that the naive estimate indicates that interaction has strong and statistically significant effects on aesthetic consumption (as given by the t-statistic).  However, once we account for unobserved components that affect both interaction and consumption---in effect, controlling for the fact that individuals who frequently engage the arts self-select into settings in which culture-talk is likely to happen---the statistical significance of the effect of interaction on consumption is greatly deflated.  In fact, only the effect of culture-mediated interaction on the frequency of reading fiction remains statistically discernible from zero at conventional levels ($t=1.98$).  The situation is very different when looking at the effect of consumption on interaction (the last four columns of table~\ref{tab3}).  Here we can see that the naive effect ($\alpha_{21}$) is an \emph{underestimate} of the true effect ($\delta_{21}$) in all domains, suggesting that the effect of consumption on interaction is deflated \emph{downwards} in the univariate probit model due to the fact that this model does not take into account the common factor that affect both consumption and interaction (in the culture conversion interpretation, this is ``embodied cultural capital''). After accounting for the confounding impact of those common factors that affect both interaction and consumption, the model shows that consumption has substantively large and statistically significant effects on the propensity to engage in culture-mediated interaction ($p<0.01$).

The estimate of the direct endogenous effect from the recursive bivariate probit model however, is not necessarily a good guide as to the total magnitude of the relationship between aesthetic consumption and interaction.  The reason for this can be seen in Figure~\ref{fig3}: in the recursive bivariate probit, each endogenous variable affects the other in two ways: directly (through the $\delta$ parameter), and \emph{indirectly} by passing on the influence of the exogenous factors common to both that are not included in the model ($\rho$).  A better estimate of the total effect of each variable is thus given by the \emph{marginal effect} \citep{long_freese06}, which is the instantaneous shift in the probability of observing high rates of social interaction given that we ``switch'' the likelihood of being a frequent consumer in that domain from zero to one (and the reverse for the recursive model predicting consumption from interaction).  The marginal effects corresponding to models (3) and (4) in Figure~\ref{fig3} are shown in Figure~\ref{fig4}.  The results are consistent with the preponderance of the culture conversion effect over the taste formation through interaction effect.  

This disparity is most evident in the case of the Movies, Dramatic Performances and the Arts.  In the case of movies, frequent consumption increases the probability of frequent use of film culture in interaction by 42\%; the corresponding marginal effect going in the oppossite direction is 0.02 (statistically indistinguishable from zero). In the realm of music, consumption increses the probability of interaction by about 40\%; while interaction appears to have comparable effect on music consumption (29\%), the z-statistic corresponding to the marginal effect estimate cannot reject the null-hypothesis that the marginal effect is in fact zero ($Z=0.45, p=0.66$).  In the case of movies, frequent attendance to galleries and museums increases the chances of culture-mediated interaction in that domain by 54\% while the corresponding marginal effect going in the oppossite direction is 0.08. 

The only realm in which the effects are of comparable size is for books. In this domain, interaction increases reading frequency and reading frequency increases the probability of talking to friends about literature by about 50\%. I believe that there is a theoretically illuminating way to account for these results. It relies on seeing the case of books as paradigmatic of the operation of an ``embedded social foci'' mechanism accounting for the connection between taste, social position and culture-talk. As noted above, books are distinctive in being the only cultural domain that has developed a widespread organizational structure in which lay consumers can organize themselves to engage in culture-mediated interaction: the book club \citep{long03}. While there exist organizations dedicated to the discussion of culture emanating from all of the cultural domains under consideration here, none has achieved the widespread institutionalization among lay audiences as the book club, having remained the purview of relatively selective (and relatively small) circles of enthusiasts \citep{kadushin66}. Thus, it is not surprising that the network model performs almost as well as the culture conversion model in the literary domain.

\subsection{The effect of intellectualized cultural dispositions on the propensity to engage in culture-mediated interaction}
Are persons who command certain forms of aesthetic orientations more likely to engage in culture-based talk?  Theory and research in culture and stratification suggests that this is likely to be the case.  Culture-talk among high-status groups is likely to be based on those qualities of aesthetic goods most likely to be the subject of standardization and institutionalization in school curricula or likely to be explicitly formulated and elaborated within informal networks of enthusiasts and cultural professionals linked to critics, journalists and other ``genre-making'' gatekeepers \citep{bourdieu93a, ostrower98, kadushin66}.  In particular, cultural orientations consistent with the institutionalized demands of highly elaborated fields of cultural production---what Bourdieu \citeyearpar{bourdieu84} referred to as ``the aesthetic disposition''---are more likely to generate culture-talk in the context of informal networks of acquaintanceship and solidarity within high-education strata \citep{collins75, dimaggio87}.  

The extent to which the aesthetic qualities of a given cultural form are more or less likely to lend themselves to verbal expression is of course not inherent in the form itself, but dependent on the institutional history of the cultural domain in question \citep{bourdieu93a}.  Drawing on recent theory and research on the phenomenon of {\em intellectulization} within artworlds \citep{baumann01,ferguson98}, we can hypothesize that the presence (or relative absence) of an already clearly specified, highly institutionalized critical vocabulary capable of being recruited for culture-based talk, may explain variation across educational categories in the extent to which we observe culture-mediated interaction in that domain.  In short, the fact that high education people are more likely to engage cultural genres that have been subject to a process of intellectualization and aesthetic elaboration by critical gatekeepers may explain their differential ability to draw on culture from that domain as a cultural resource in conservation, even though they are no more likely to \emph{participate} in that broad domain than low education respondents. 

For instance, as we saw above, members of all educational groups are about equally likely to report listening to music very often; however, high education respondents are substantially more likely to report {\em talking} about music with their friends than their low education counterparts. If the above considerations are correct, this phenomenon would be explainable by the fact that there exists variation in culture-based talk in the music domain across persons who belong to different \emph{taste cultures} \citep{gans99}.  In particular, respondents more likely to display a preference for those genre categories more likely to attract high-status individuals should be more likely to generate culture-mediated interaction than persons who display a preference for those genres more likely to attract individuals of lower status. Thus, heterogeneity among music listeners in terms of allegiance to different taste cultures may partially account for why we observe educational heterogeneity in terms of the deployment of musical culture as a resource in interaction without a corresponding form of education-based stratification in terms of the frequency of listening to music. 

\subsection{Multivariate probit specification}               
To explore this possibility, I use information on the respondent's favorite musical genres to explore the effect of aesthetic orientation---high-status versus ``popular''---on the propensity to engage in culture-based talk.\footnote{Most respondents have only one favorite. Only a small minority chooses more than two.}.  The ``intellectualization'' hypothesis suggests that it is precisely those persons who have developed an aesthetic disposition toward those genres most clearly institutionalized as ``art'' that will be more likely to engage in culture-based talk (net of the effect of consumption frequency).  I specify a multivariate probit model with correlated errors of the form shown in Figure~\ref{fig4}.  The model includes an additional endogenous covariate (favorite musical genre; $y_{3}$) that is affected by a vector of socio-demographic characteristics $\beta_{3}x$, and that is also affected by consumption frequency ($y_{3}$). Aesthetic orientation is allowed to affect frequency of culture-based talk involving music via the $\delta_{31}$ pathway.  Because (as we have seen) talk those not affect consumption, we can set $\delta{12}$ to zero and the trivariate probit model is identified. The model accounts for those exogenous factors that affect both taste for high-status culture and culture-based talk ($\rho_{13}$), thus allowing to cleanly identify the effect that goes from taste to culture-mediated interaction. We should expect this effect to be positive for those respondents reporting a preference for high-status genres, and null (or negative) for the rest.

The results of the multivariate probit model are shown in Table~\ref{tab6}.  Although the fact that reports on taste are forced to conform to this highly aggregate (and somewhat heterogeneous) genre categories shown in the table is not ideal, the overall pattern of results conforms to expectations. Taste for Classical music or Opera has a positive and statistically significant effect on the rates of culture-based talk involving music ($t=3.03$).  The effect of liking Jazz or Blues is in the expected direction, but is not statistically significant.  For the lower status genres, the effect approaches zero and becomes negative for Christian and Gospel and Punk and Heavy Metal.  The results are thus consistent with the hypothesis that command of the codes that govern those musical genres most likely to have already institutionalized organizational and critical resources for intellectualization, produces culture-based talk. Respondents who like those genres most associated with either ``scenes'' or ``industry'' \citep{lena_peterson08} on the other hand are no more likely to generate conversation than expected by chance after controlling for the respondent's socio-demographic profile and for the effect of consumption frequency.  These results help to partly explain why we observe stratification of culture-mediated interaction in the music domain in the absence of education-linked consumption differences.

Table~\ref{tab7} shows evidence that is consistent with the claim that an orientation towards high-status cultural forms generates culture-talk \emph{across} domains. In the table I show the coefficient estimates corresponding to the same multivariate probit model shown in Figure~\ref{fig5}, but with culture-based talk in the domains of Movies, Literature, Arts and Dramatic Performances playing the role of $y_{1}$.  The results show that orientation toward high status culture---measured by the same variable as before; namely liking classical or opera---has a positive effect on culture-based talk in all domains ($p<0.01$) with the exception of literature ($t=-0.44$), net of the effect of frequency of consumption in that domain and the other socio-demographic predictors. 

\section{Discussion}
\subsection{Summary of the Results}
The goals of these paper have been twofold: first to highlight the theoretical and empirical importance of the phenomenon that I have referred to as {\em culture-mediated interaction}. Beginning with the Simmelian legacy on the importance of seemingly ``trivial'' conversational content for the emergence of ``sociability'' (or the ``play form'' of association) to the contemporary stress on ``entertainment talk'' as the cement that holds status groups together as micro-relational communities of acquaintance \citep{collins75, dimaggio87}, I have argued that the deployment of cultural knowledge in concrete interaction is one of the central, but least understood, processes in the field of culture stratification \citep{purcell07}.  In spite of its centrality, it has been a relatively under-studied phenomenon. This paper has attempted to redress this issue, but examining the socio-demographic distribution and cultural correlates of the tendency to talk to friends about culture consumption experiences.  

The second goal of this paper was to advance research on the subject of the role of culture-based talk as a constitutive component of lifestyle-based stratification in contemporary post-traditional societies, as well as to clarify the relationship between aesthetic consumption, culture-talk and cultural taste. The results show that culture-based talk is a relatively common experience for the large majority of Americans, with almost four fifths reporting engaging in some form of culture-mediated interaction using cultural knowledge from the five domains included in the survey. The results also show however, (consistent with prior theory in cultural stratification) that culture-mediated interaction is stratified by the same markers of social position (Age, Gender) and social status (Education) that previous research has shown also characterize audience segmentation in the corresponding cultural domains.  

In particular, culture-talk in those domains characterized by large-scale commercial culture (Movies, Music) is strongly stratified by Age (with the young being more likely to use such culture in conversation) and Education (with the highly educated generating more culture-based talk). Culture-talk recruiting cultural knowledge from the higher-status ``non-profit'' arts is largely stratified by educational attainment.  Literature is distinctive in that in addition to exhibiting the expected substantial rates of stratification by education it is also strongly stratified by gender (although all domains except for music show a female advantage in culture-based talk) as suggested by recent research on gender differences in recruitment into the ``reading class.''  

Nevertheless, it is important to note that relative rates of culture talk increase with education in all domains, whether commercial, literary or artistic.  The educational stratum most likely to recruit cultural knowledge in interaction consists of those who have been bestowed a postgraduate degree. This is consistent with the contention that culture-talk is an interactional resource characteristic of high-status groups in contemporary American society. This is also consistent with prior research that shows that high-aesthetic consumption groups, are also those more likely to possess richer networks of strong ties (which provide closure).  Such markers of social position as marital status, race and ethnicity and urban residence in contrast, are distinctive for the fact that they do not serve as differentiating factors in predicting the propensity of culture-mediated interaction.

This paper also explored an important question in the sociological study of culture taste: the interrelationship between propensity to consumer culture, rates of culture-mediated interaction and markers of social position.  Applying models for the determination of directional effect in the presence of ``endogeneity'' I showed that a model that conceptualizes rates of aesthetic consumption as the primary driver for the generation of opportunities to produce culture-based talk outperforms a model that sees interaction opportunities as the generator of tastes for certain consumption derived from network theory.  In all domains (and consistent with a view of tastes as relatively durable aspects of the person's embodied cultural capital), frequency of consumption generate higher propensities to use culture in interaction, but not the reverse (with the exception of literature, which exhibits bi-directional effects).  

While for the most part the stratification of culture-based talk in each domain mirrors the stratification of aesthetic consumption in that same consumption realm, important discrepancies are evident. In particular, the realms more deeply connected to ``commercial'' culture (movies and music) exhibit much weaker stratification by education in terms of consumption than they do in terms of talk.  This suggests that high-status individuals are able to generate culture-mediated interaction using knowledge procured from so-called ``popular'' cultural domains even as their consumption rates in these domains do not differ as dramatically from the less educated as they do for the high status domains (literature and the non-profit arts).  

Integrating aesthetic orientations towards high status culture into a model predicting culture-based talk propensities, I showed that a model that incorporates recent research and theory on ``intellectualization'' in cultural domains institutionalized as artworlds, generates the testable prediction that it is those persons more likely to be oriented toward high-status culture that will be more likely to generate culture-based talk within the domains putatively dedicated to the commercial arts. The results show that persons exhibiting high-status taste patterns in the domain of music, are more likely to generate culture-based talk (net of the effect of consumption frequency) within that domain but also across other domains. This suggests that culture-mediated interaction among high-status classes may be produced by the ``transposable'' aptitude that Bourdieu \citeyearpar{bourdieu84} referred to as ``the aesthetic disposition.'' 
                
               
\subsection{Limitations of the analysis}
Naturally, the conclusions outlined above should be taken as preliminary suggestions and guidelines for future research rather than as firmly established propositions.  Obviously the most clear limitation of this study concerns the limited range of cultural domain included in the NSS for which information on the rate of culture-mediated interaction is available.  In particular, as befits an arts participation survey, questions about culture-mediated interaction are weighted toward the most legitimate forms of aesthetic practice but are silent when it comes to more popular culture consumption experiences \citep{dimaggio_mukhtar04}. 
                
In particular, past research suggests that culture-mediated interaction centered around sports \citep{erickson96}; Fashion and Cuisine \citep{dimaggio87} and in particular television \citep{fiske87} are a consistent generator of culture-based talk in formal  and informal settings (e.g. the workplace).  Thus, we can be certain that the rates of culture based talk reported above underestimate by a substantial amount its prevalence in the American population.  
                
The absence of information on culture-based talk in more popular domains may also lead to an overstatement of its connection to markers of stratification, such as educational credentials; although it is reasonable to suppose that just like in the case of music and movies, we may be likely to observe education-based stratification in rates of culture-mediated interaction in the absence of education-based stratification in consumption rates.  
                
It is also likely that including a wider range of cultural practices and domains may heavily qualify the finding that it is those more legitimate cultural pursuits within domains that generate the most culture-based talk.  It is likely for instance that within the realm of television-programming, it is those genres endowed with the least amount of critical legitimacy (e.g. reality television) that generate the most culture-based talk, especially among the young.   The very limited research done on ``sports talk'' \citep{purcell07} suggests however, that ``intellectualization'' in the realm of sports (e.g. produced by adopting the critical vocabulary produced by sports journalists) does facilitate sports-based talk among the educated, as suggested by the above analysis. 
                
The other major limitation of the present study is its relative glossing of the {\em context} of consumption as a mediator of the extent to which given aesthetic consumption experiences will generate culture-based talk. Some genres have built-in rituals of audience reception in which ``interaction'' is expected to occur {\em within} the consumption site and {\em during} the performance (e.g. a rock concert) and which thus carry little expectations that the cultural-knowledge gained therein will be transposed to other informal contexts as a generator of conversation.  Culture that is consumed according to socially constructed standards of audience ``propriety'' (e.g. sitting quietly in rapt attention) may in turn generate a lot of culture-based talk precisely because direct interaction among audience members (verbal and non-verbal) is frowned upon during the process of reception.  Cultural domains that exhibit wide variation in the context (and thus the expected audience conventions) of reception, may be seen to have a ``bi-modal'' distribution of their likelihood of being used for culture-based talk (like music in these data), with non-negligible proportions of active consumers reporting not using (as well as using) culture for that domain as a conversational resource.                
                
                   
\section{Conclusion: why cultural matters may matter more than ``important'' matters}
The pervasiveness and social distribution of the phenomenon of culture-mediated interaction has important implications for the ways in which data about social connectivity are collected and conceptualized.  The two large-scale efforts to collect data about patterns of network connectivity in the American population have relied on the ``important matters'' name generator methodological tool \citep{burt84, marsden87, mcpherson_etal06}. The presumption here is that there is some sort of homology between the alleged {\em importance of the content of conversation and the importance of the tie in question}.  Importance here being conceptualized both as 1) theoretical importance for the social scientist,--e.g. the status of the social ties identified in this manner as an measure of the extent to which the individual is ``connected'' \citep{burt90}; and 2) practical importance for the person; or the extent to which the connection identified in this manner give us access to a meaningful subset of the relations that a person would draw on in time of emotional or financial help \citep{wellman_wortley90}.
                  
But given the social distribution of culture -mediated interaction---and its intuitive connection with certain socio-demographic markers---it is perilous to suppose that the ``importance'' of the ``matters'' discussed has any connection to the ``importance'' of the contact with which they are discussed.  In fact, given the theoretical framework outlined here, we should expect little or no connection between the objective importance of the network content and the object (or social-scientific) importance of the ties connected to to those contents.  If Simmelian ``pure sociability'' occurs within those relationships which provide {\em the most intrinsic satisfaction} we should expect the connection to even be negative (e.g. meaningful and subjectively important ties are sustained through seemingly ``trivial'' and ``unimportant'' cultural content).  This becomes more important when we realize that, as Simmel would have predicted, conversation involving ``important matters'' is too sporadic and inconsistent a phenomenon with which to judge whether persons are connected or not \citep{bearman_parigi04}.
                  
A substantial number of people (including a larger proportion of married persons) report simply not having anybody with which to discuss ``important'' matters with \citep{bearman_parigi04}). Should we conclude therefore that most persons are ``disconnected'' and experience little sociability \citep{mcpherson_etal06, putnam00}?  I would suggest that a more theoretically sophisticated understanding of what ``sociability'' is all about suggests otherwise. ``Important matters'' name generators may not be getting at either ``important'' content  or important ties (if the social-scientific interest is to uncover those ties that are meaningful for the individual).  This becomes even more of an issue when we consider that interaction based on important matters, may be more steeply stratified by markers of social status, such as education, than culture-based talk \citep{bourdieu84, mcpherson_etal06, marsden87, fischer82}.
                  
In this respect, the ``important matters'' name-generator appears to be a tough ``test'', that only a few privileged persons endowed with certain forms of linguistic (or political) capital may pass \citep{bourdieu84}. Important matters name generators may thus confound stratification based on aptitudes orthogonal to social competence and ability to form connections with stratification criteria that directly impinge on the persons' ability to in fact create and sustain social relationships \citep{lin99, dimaggio96, lizardo06a}
                  
If the framework outlined here is on the right track, {\em cultural matters may matter more than important matters} as the content with which to sustain (and thus measure) meaningful ties.  Thus, querying persons as to whom they discuss the ``trivial'' matters related to the cultural content that makes up entertainment talk (e.g. sports, fashion, cuisine, television, movies, the arts), maybe a more sensible strategy to generate data on personal social connectivity, and to  draw conclusions about the fate of this connectivity in the contemporary situation.

\newpage
\bibliographystyle{asr}
\bibliography{N:/Private/zotero-library}

\newpage
\begin{landscape}
\begin{table}[htb!]
\begin{tabular*}{\hsize}{@{\extracolsep{\fill}} l*{7}{c}} 
\hline
\hline
&\multicolumn{2}{c}{Movies}&\multicolumn{2}{c}{Music}&\multicolumn{2}{c}{Books}\\
\hline
&Talk&Consumption&Talk&Consumption&Talk&Consumption\\
Education (Less than H.S.=1)&-0.37 (-5.05)&-0.42 (-5.56)&-0.37 (-4.94)&-0.04 (-0.49)&-0.39 (-5.20)&-0.40 (-5.45)\\
Education (High School=1)&-0.14 (-2.80)&-0.11 (-2.14)&-0.24 (-4.96)&-0.09 (-1.71)&-0.29 (-6.00)&-0.24 (-4.96)\\
Education (Some Coll.=1)&0.06 (1.19)&0.14 (2.79)&0.11 (2.31)&0.03 (0.51)&0.12 (2.49)&0.07 (1.44)\\
Education (College=1)&0.11 (2.00)&0.13 (2.25)&0.17 (3.25)&-0.01 (-0.17)&0.25 (4.79)&0.20 (3.59)\\
Education (Postgraduate=1)&0.34 (4.51)&0.26 (3.73)&0.33 (4.76)&0.11 (1.37)&0.31 (4.36)&0.36 (5.26)\\
Age (18-24=1)   &0.51 (6.57)&0.68 (8.48)&0.39 (5.16)&0.59 (6.20)&0.18 (2.33)&0.03 (0.41)\\
Age (25-29=1)   &0.37 (4.56)&0.37 (4.43)&0.21 (2.74)&0.16 (1.64)&0.07 (0.87)&-0.09 (-1.08)\\
Age (30-34=1)   &0.07 (0.91)&-0.08 (-1.05)&-0.04 (-0.55)&0.03 (0.31)&0.02 (0.29)&-0.10 (-1.22)\\
Age (35-39=1)   &0.15 (1.82)&0.05 (0.63)&-0.00 (-0.03)&0.06 (0.59)&-0.07 (-0.84)&-0.12 (-1.52)\\
Age (40-49=1)   &-0.01 (-0.15)&-0.13 (-2.09)&-0.08 (-1.29)&-0.15 (-2.20)&-0.01 (-0.21)&0.01 (0.17)\\
Age (50-64=1)   &-0.30 (-5.34)&-0.24 (-4.33)&-0.13 (-2.31)&-0.22 (-3.52)&-0.01 (-0.24)&0.08 (1.48)\\
Age (65+=1)     &-0.80 (-12.38)&-0.65 (-10.00)&-0.34 (-5.46)&-0.46 (-6.85)&-0.18 (-2.98)&0.19 (3.14)\\
Income (75K+=1) &0.09 (2.27)&0.18 (4.44)&0.07 (1.79)&0.04 (0.96)&0.01 (0.24)&0.02 (0.44)\\
Gender (Women=1)&0.13 (4.95)&0.07 (2.44)&0.05 (1.72)&0.08 (2.60)&0.22 (8.15)&0.17 (6.44)\\
Race/Eth. (Black/Afr. Amer.=1)&-0.02 (-0.48)&0.03 (0.71)&0.03 (0.65)&0.12 (2.25)&0.06 (1.34)&0.09 (1.97)\\
Race/Eth. (Hispanic=1)&-0.03 (-0.56)&0.17 (2.84)&-0.07 (-1.09)&-0.06 (-0.75)&-0.03 (-0.53)&0.12 (1.90)\\
Mar. Status (Married=1)&-0.02 (-0.85)&-0.02 (-0.79)&-0.05 (-1.63)&-0.04 (-1.27)&-0.02 (-0.63)&-0.04 (-1.30)\\
Children (Yes=1)&0.04 (1.21)&0.06 (1.76)&-0.00 (-0.12)&0.09 (2.05)&0.01 (0.38)&-0.01 (-0.28)\\
Urban Residence (Yes=1)&0.04 (1.42)&0.09 (3.20)&0.06 (2.34)&0.01 (0.47)&0.03 (0.96)&0.00 (0.06)\\
Year of Survey (2004=1)&-0.03 (-1.13)&-0.13 (-4.69)&-0.01 (-0.25)&-0.12 (-4.10)&-0.01 (-0.37)&0.01 (0.33)\\
Observations    &     1729&         &     1733&         &     1731&         \\
\hline
\end{tabular*}
\caption{}
\label{tab1}
\end{table}
\end{landscape}

\begin{landscape}
\begin{table}[htb!]
\begin{tabular*}{\hsize}{@{\extracolsep{\fill}} l*{5}{c}} 
\hline
\hline
&\multicolumn{2}{c}{Arts}&\multicolumn{2}{c}{Plays/Drama}\\
\hline
&Talk&Consumption&Talk&Consumption\\
Education (Less than H.S.=1)&-0.23 (-2.73)&-0.47 (-5.54)&-0.28 (-3.46)&-0.47 (-6.17)\\
Education (High School=1)&-0.40 (-7.57)&-0.41 (-7.45)&-0.34 (-6.84)&-0.27 (-5.47)\\
Education (Some Coll.=1)&0.02 (0.30)&0.09 (1.73)&0.05 (1.01)&0.10 (2.12)\\
Education (College=1)&0.12 (2.13)&0.23 (4.09)&0.15 (2.82)&0.20 (3.72)\\
Education (Postgraduate=1)&0.49 (7.02)&0.55 (7.76)&0.42 (5.97)&0.43 (6.31)\\
Age (18-24=1)   &0.05 (0.61)&0.09 (1.09)&0.16 (2.18)&0.25 (3.28)\\
Age (25-29=1)   &0.07 (0.88)&0.08 (0.97)&0.09 (1.16)&0.15 (1.73)\\
Age (30-34=1)   &0.05 (0.66)&0.06 (0.72)&0.06 (0.67)&-0.04 (-0.45)\\
Age (35-39=1)   &0.11 (1.36)&-0.12 (-1.34)&0.02 (0.25)&-0.08 (-1.03)\\
Age (40-49=1)   &-0.01 (-0.10)&-0.00 (-0.04)&0.05 (0.80)&-0.03 (-0.41)\\
Age (50-64=1)   &-0.07 (-1.10)&-0.02 (-0.32)&-0.14 (-2.47)&-0.12 (-2.16)\\
Age (65+=1)     &-0.21 (-3.12)&-0.09 (-1.35)&-0.25 (-3.80)&-0.13 (-2.04)\\
Income (75K+=1) &0.05 (1.32)&0.14 (3.50)&0.13 (3.46)&0.13 (3.16)\\
Gender (Women=1)&0.11 (4.03)&0.05 (1.69)&0.14 (4.98)&0.02 (0.58)\\
Race/Eth. (Black/Afr. Amer.=1)&0.03 (0.62)&-0.02 (-0.37)&0.11 (2.59)&0.01 (0.26)\\
Race/Eth. (Hispanic=1)&0.05 (0.74)&0.00 (0.02)&-0.02 (-0.29)&0.03 (0.51)\\
Mar. Status (Married=1)&-0.04 (-1.32)&-0.04 (-1.42)&-0.02 (-0.85)&-0.05 (-1.54)\\
Children (Yes=1)&-0.06 (-1.51)&-0.02 (-0.44)&-0.05 (-1.23)&-0.01 (-0.22)\\
Urban Residence (Yes=1)&0.05 (1.73)&0.04 (1.52)&0.08 (3.07)&0.07 (2.42)\\
Year of Survey (2004=1)&-0.05 (-1.71)&-0.03 (-0.96)&-0.05 (-1.93)&-0.12 (-4.46)\\
Observations    &     1728&         &     1731&         \\
\hline
\end{tabular*}
\caption{}
\label{tab2}
\end{table}
\end{landscape}


\pagebreak
\begin{landscape}
\begin{table}[htb!]
\begin{tabular*}{\hsize}{@{\extracolsep{\fill}} l*{13}{c}} 
\hline 
\hline           
\multicolumn{13}{c}{Culture-Mediated Interaction}\\
&\multicolumn{2}{c}{Education (\emph{df}=4)}&\multicolumn{2}{c}{Age (\emph{df}=6)}&\multicolumn{2}{c}{Gender (\emph{df}=1)}&\multicolumn{2}{c}{Race, Eth. (\emph{df}=2)}&\multicolumn{4}{c}{Model}\\
& $\chi^{2}$  & {\em p} & $\chi^{2}$ & {\em p} &  $\chi^{2}$ &  {\em p} &  $\chi^{2}$ &  {\em p} &Educ\% &  Age\%  &  Gender\% & Race\%\\ 
\hline
Movies&43.3&0.000&203.9&0.000&24.5&0.000&0.5&0.785&0.148&0.694&0.083&0.002\\
Music&68.4&0.000&61.6&0.000&3.0&0.085&1.8&0.414&0.412&0.371&0.018&0.011\\
Books&88.0&0.000&14.6&0.024&66.5&0.000&2.3&0.320&0.517&0.085&0.390&0.013\\
Arts&95.2&0.000&11.5&0.075&16.3&0.000&0.8&0.655&0.656&0.079&0.112&0.006\\
Plays&83.9&0.000&23.3&0.001&24.8&0.000&7.1&0.029&0.479&0.133&0.142&0.040\\
\hline \\ \\
\hline
\hline
\multicolumn{13}{c}{Aesthetic Consumption}\\
&\multicolumn{2}{c}{Education (\emph{df}=4)}&\multicolumn{2}{c}{Age (\emph{df}=6)}&\multicolumn{2}{c}{Gender (\emph{df}=1)}&\multicolumn{2}{c}{Race, Eth. (\emph{df}=2)}&\multicolumn{4}{c}{Model}\\
& $\chi^{2}$  & {\em p} & $\chi^{2}$ & {\em p} &  $\chi^{2}$ &  {\em p} &  $\chi^{2}$ &  {\em p} &Educ\% &  Age\%  &  Gender\% & Race\%\\
\hline
Movies&46.6&0.000&176.0&0.000&5.9&0.015&8.2&0.016&0.142&0.536&0.018&0.025\\
Music&4.3&0.369&84.3&0.000&6.8&0.009&6.1&0.048&0.031&0.601&0.048&0.043\\
Books&74.1&0.000&13.0&0.043&41.5&0.000&6.7&0.035&0.559&0.098&0.313&0.050\\
Arts&131.1&0.000&5.5&0.480&2.9&0.091&0.1&0.931&0.725&0.030&0.016&0.001\\
Plays&98.7&0.000&19.7&0.003&0.3&0.565&0.3&0.861&0.562&0.112&0.002&0.002\\
\hline
\end{tabular*}
\caption{}
\label{tab3}
\end{table}
\end{landscape}

\pagebreak
\begin{landscape}
\begin{table}[htb!]
\begin{tabular*}{\hsize}{@{\extracolsep{\fill}} l*{9}{c}} 
\hline 
\hline 
&\multicolumn{2}{c}{Education}&\multicolumn{2}{c}{Age}&\multicolumn{2}{c}{Gender}&\multicolumn{2}{c}{Race/Eth.}\\
&\multicolumn{2}{c}{$H_{0}: \beta_{1}=\beta_{2}$}&\multicolumn{2}{c}{$H_{0}: \beta_{1}=\beta_{2}$}&\multicolumn{2}{c}{$H_{0}: \beta_{1}=\beta_{2}$}&\multicolumn{2}{c}{$H_{0}:  \beta_{1}=\beta_{2}$}\\
&$\chi^{2}$&p-value&$\chi^{2}$&p-value&$\chi^{2}$&p-value&$\chi^{2}$&p-value\\
\hline 
Movies&3.03&0.553&12.77&0.047&5.31&0.021&8.54&0.014\\
Music&24.32&0.000&7.64&0.265&0.94&0.331&2.12&0.346\\
Books&3.44&0.487&36.60&0.000&2.44&0.118&5.15&0.076\\
Arts&9.45&0.051&9.36&0.154&5.39&0.020&1.21&0.545\\
Plays&4.90&0.298&6.51&0.368&17.55&0.000&5.09&0.079\\
\hline 
\end{tabular*}
\caption{}
\label{tab4}
\end{table}
\end{landscape}

\pagebreak
\begin{landscape}
\begin{table}[htb!]                                 
\begin{tabular*}{\hsize}{@{\extracolsep{\fill}} l*{9}{c}} 
\hline 
\hline
&\multicolumn{2}{c}{Baseline}&\multicolumn{2}{c}{SUR}&\multicolumn{2}{c}{Model 3}&\multicolumn{2}{c}{Model 4} \\
&$\rho_{12}$&\emph{p}-value&$\rho_{12}$&\emph{p}-value&$\rho_{12}$&\emph{p}-value&$\rho_{12}$&\emph{p}-value\\
\hline 
Movies&0.64&0.00&0.60&0.00&-0.35&0.50&-0.70&0.22\\
Music&0.40&0.00&0.37&0.00&0.77&0.59&-0.83&0.01\\
Books&0.63&0.00&0.60&0.00&-0.41&0.72&-0.25&0.64\\
Arts&0.72&0.00&0.70&0.00&0.26&0.80&-0.24&0.69\\
Plays&0.57&0.00&0.55&0.00&-0.52&0.23&0.87&0.04\\
\hline \\
\hline 
\hline
&\multicolumn{2}{c}{Model 1}&\multicolumn{2}{c}{Model 3}&\multicolumn{2}{c}{Model 2}&\multicolumn{2}{c}{Model 4} \\
&$\alpha_{12}$&\emph{t}-stat&$\delta_{12}$&\emph{t}-stat&$\alpha_{21}$&\emph{t}-stat&$\delta_{21}$&\emph{t}-stat \\
\hline
Movies&1.15&13.04&1.67&1.69&1.15&12.94&2.11&8.72\\
Music&0.62&9.27&-0.72&-1.01&0.64&9.43&1.77&11.22\\
Books&1.09&16.28&1.69&1.25&1.09&16.34&1.48&2.08\\
Arts&1.40&15.78&0.92&0.42&1.44&15.74&1.86&2.37\\
Plays&1.01&13.49&1.81&4.23&0.99&13.32&-0.65&-1.86\\
\hline
\end{tabular*}
\caption{}
\label{tab5}
\end{table}
\end{landscape}

\pagebreak
\begin{landscape}
\begin{table}[htb!]                                 
\begin{tabular*}{\hsize}{@{\extracolsep{\fill}} l*{5}{c}} 
\hline 
\hline
&\multicolumn{4}{c}{Movies}\\
&$pr(y_{1}=0 \wedge y_{2}=1)$&$pr(y_{1}=1 \wedge y_{2}=0)$&$pr(y_{1}=1 \wedge y_{2}=1)$&$pr(y_{1}=0 \wedge y_{2}=0)$\\
\hline
Education (Less than H.S.=1) (d)&   0.00&  -0.06&  -0.07&   0.13\\
Education (High School=1) (d)&   0.01&  -0.03&  -0.02&   0.05\\
Education (Some Coll.=1) (d)&   0.00&  -0.01&   0.02&  -0.01\\
Education (College=1) (d)&  -0.00&   0.02&   0.03&  -0.05\\
Education (Postgraduate=1) (d)&  -0.01&   0.07&   0.06&  -0.11\\
Age (18-24=1) (d)&   0.01&  -0.05&   0.23&  -0.18\\
Age (25-29=1) (d)&  -0.01&   0.03&   0.11&  -0.13\\
Age (30-34=1) (d)&  -0.01&   0.04&  -0.01&  -0.02\\
Age (35-39=1) (d)&  -0.01&   0.04&   0.03&  -0.07\\
Age (40-49=1) (d)&  -0.01&   0.03&  -0.03&   0.00\\
Age (50-64=1) (d)&   0.01&  -0.06&  -0.07&   0.12\\
Age (65+=1) (d) &   0.00&  -0.15&  -0.19&   0.33\\
Income (75K+=1) (d)&   0.00&  -0.00&   0.04&  -0.04\\
Gender (Women=1) (d)&  -0.01&   0.03&   0.03&  -0.05\\
Race/Eth. (Black/Afr. Amer.=1) (d)&   0.01&  -0.03&   0.03&  -0.00\\
Race/Eth. (Hispanic=1) (d)&   0.02&  -0.07&   0.04&   0.01\\
\hline
\end{tabular*}
\caption{}
\label{tab6a}
\end{table}
\end{landscape}

\pagebreak
\begin{landscape}
\begin{table}[htb!]                                 
\begin{tabular*}{\hsize}{@{\extracolsep{\fill}} l*{5}{c}} 
\hline 
\hline
&\multicolumn{4}{c}{Music}\\
&$pr(y_{1}=0 \wedge y_{2}=1)$&$pr(y_{1}=1 \wedge y_{2}=0)$&$pr(y_{1}=1 \wedge y_{2}=1)$&$pr(y_{1}=0 \wedge y_{2}=0)$\\
\hline
Education (Less than H.S.=1) (d)&   0.08&  -0.04&  -0.08&   0.04\\
Education (High School=1) (d)&   0.05&  -0.02&  -0.08&   0.05\\
Education (Some Coll.=1) (d)&  -0.03&   0.02&   0.03&  -0.01\\
Education (College=1) (d)&  -0.03&   0.02&   0.03&  -0.02\\
Education (Postgraduate=1) (d)&  -0.07&   0.03&   0.10&  -0.06\\
Age (18-24=1) (d)&  -0.00&  -0.04&   0.21&  -0.16\\
Age (25-29=1) (d)&   0.01&  -0.02&   0.07&  -0.06\\
Age (30-34=1) (d)&   0.04&  -0.02&  -0.02&   0.00\\
Age (35-39=1) (d)&  -0.00&   0.00&  -0.00&   0.00\\
Age (40-49=1) (d)&  -0.01&   0.01&  -0.05&   0.05\\
Age (50-64=1) (d)&  -0.03&   0.02&  -0.06&   0.06\\
Age (65+=1) (d) &  -0.05&   0.04&  -0.14&   0.15\\
Income (75K+=1) (d)&  -0.02&   0.01&   0.02&  -0.01\\
Gender (Women=1) (d)&   0.01&  -0.01&   0.03&  -0.03\\
Race/Eth. (Black/Afr. Amer.=1) (d)&   0.01&  -0.01&   0.05&  -0.05\\
Race/Eth. (Hispanic=1) (d)&  -0.01&   0.01&  -0.02&   0.02\\
\hline
\end{tabular*}
\caption{}
\label{tab6b}
\end{table}
\end{landscape}

\pagebreak
\begin{landscape}
\begin{table}[htb!]                                 
\begin{tabular*}{\hsize}{@{\extracolsep{\fill}} l*{5}{c}} 
\hline 
\hline
&\multicolumn{4}{c}{Books}\\
&$pr(y_{1}=0 \wedge y_{2}=1)$&$pr(y_{1}=1 \wedge y_{2}=0)$&$pr(y_{1}=1 \wedge y_{2}=1)$&$pr(y_{1}=0 \wedge y_{2}=0)$\\
\hline

Education (Less than H.S.=1) (d)&  -0.00&  -0.00&  -0.15&   0.15\\
Education (High School=1) (d)&   0.03&  -0.02&  -0.12&   0.11\\
Education (Some Coll.=1) (d)&  -0.00&  -0.00&   0.05&  -0.04\\
Education (College=1) (d)&  -0.02&   0.00&   0.10&  -0.09\\
Education (Postgraduate=1) (d)&  -0.03&  -0.00&   0.14&  -0.11\\
Age (18-24=1) (d)&  -0.01&   0.01&   0.05&  -0.05\\
Age (25-29=1) (d)&  -0.03&   0.02&  -0.00&   0.00\\
Age (30-34=1) (d)&  -0.03&   0.02&   0.00&   0.00\\
Age (35-39=1) (d)&  -0.01&   0.01&  -0.05&   0.05\\
Age (40-49=1) (d)&   0.02&  -0.01&  -0.01&   0.01\\
Age (50-64=1) (d)&  -0.01&   0.00&   0.02&  -0.01\\
Age (65+=1) (d) &   0.07&  -0.05&  -0.01&  -0.01\\
Income (75K+=1) (d)&   0.01&  -0.01&   0.00&  -0.00\\
Gender (Women=1) (d)&  -0.01&   0.00&   0.08&  -0.07\\
Race/Eth. (Black/Afr. Amer.=1) (d)&  -0.01&   0.00&   0.03&  -0.02\\
Race/Eth. (Hispanic=1) (d)&   0.04&  -0.03&   0.02&  -0.03\\
\hline
\end{tabular*}
\caption{}
\label{tab6c}
\end{table}
\end{landscape}

\pagebreak
\begin{landscape}
\begin{table}[htb!]                                 
\begin{tabular*}{\hsize}{@{\extracolsep{\fill}} l*{5}{c}} 
\hline 
\hline
&\multicolumn{4}{c}{Arts}\\
&$pr(y_{1}=0 \wedge y_{2}=1)$&$pr(y_{1}=1 \wedge y_{2}=0)$&$pr(y_{1}=1 \wedge y_{2}=1)$&$pr(y_{1}=0 \wedge y_{2}=0)$\\
\hline
Education (Less than H.S.=1) (d)&  -0.03&  -0.01&  -0.04&   0.08\\
Education (High School=1) (d)&  -0.02&  -0.04&  -0.06&   0.12\\
Education (Some Coll.=1) (d)&   0.00&   0.01&   0.01&  -0.02\\
Education (College=1) (d)&   0.02&  -0.00&   0.03&  -0.05\\
Education (Postgraduate=1) (d)&   0.03&   0.04&   0.11&  -0.18\\
Age (18-24=1) (d)&   0.03&  -0.02&   0.01&  -0.02\\
Age (25-29=1) (d)&   0.01&  -0.01&   0.00&  -0.00\\
Age (30-34=1) (d)&  -0.00&   0.01&   0.01&  -0.02\\
Age (35-39=1) (d)&  -0.03&   0.05&  -0.01&  -0.01\\
Age (40-49=1) (d)&  -0.00&   0.01&   0.00&  -0.01\\
Age (50-64=1) (d)&   0.00&   0.00&   0.00&  -0.01\\
Age (65+=1) (d) &   0.00&  -0.03&  -0.02&   0.05\\
Income (75K+=1) (d)&   0.01&  -0.01&   0.01&  -0.01\\
Gender (Women=1) (d)&  -0.00&   0.02&   0.01&  -0.02\\
Race/Eth. (Black/Afr. Amer.=1) (d)&   0.00&   0.01&   0.01&  -0.03\\
Race/Eth. (Hispanic=1) (d)&  -0.01&   0.03&   0.01&  -0.03\\
\hline
\end{tabular*}
\caption{}
\label{tab6d}
\end{table}
\end{landscape}

\pagebreak
\begin{landscape}
\begin{table}[htb!]                                 
\begin{tabular*}{\hsize}{@{\extracolsep{\fill}} l*{5}{c}} 
\hline 
\hline
&\multicolumn{4}{c}{Plays}\\
&$pr(y_{1}=0 \wedge y_{2}=1)$&$pr(y_{1}=1 \wedge y_{2}=0)$&$pr(y_{1}=1 \wedge y_{2}=1)$&$pr(y_{1}=0 \wedge y_{2}=0)$\\
\hline
Education (Less than H.S.=1) (d)&  -0.05&  -0.02&  -0.08&   0.15\\
Education (High School=1) (d)&  -0.02&  -0.03&  -0.06&   0.11\\
Education (Some Coll.=1) (d)&   0.01&   0.01&   0.02&  -0.03\\
Education (College=1) (d)&   0.02&   0.01&   0.04&  -0.06\\
Education (Postgraduate=1) (d)&   0.04&   0.02&   0.12&  -0.18\\
Age (18-24=1) (d)&   0.07&  -0.03&   0.04&  -0.08\\
Age (25-29=1) (d)&   0.02&  -0.00&   0.02&  -0.03\\
Age (30-34=1) (d)&  -0.02&   0.02&   0.00&  -0.00\\
Age (35-39=1) (d)&  -0.02&   0.00&  -0.03&   0.05\\
Age (40-49=1) (d)&  -0.02&   0.02&   0.00&  -0.00\\
Age (50-64=1) (d)&  -0.02&   0.01&  -0.01&   0.02\\
Age (65+=1) (d) &   0.01&  -0.02&  -0.02&   0.03\\
Income (75K+=1) (d)&   0.01&   0.01&   0.03&  -0.05\\
Gender (Women=1) (d)&  -0.01&   0.02&   0.02&  -0.03\\
Race/Eth. (Black/Afr. Amer.=1) (d)&  -0.01&   0.03&   0.04&  -0.06\\
Race/Eth. (Hispanic=1) (d)&   0.01&   0.00&   0.02&  -0.03\\
\hline
\end{tabular*}
\caption{}
\label{tab6e}
\end{table}
\end{landscape}






\begin{landscape}
\begin{table}[htb!] 
\begin{tabular*}{\hsize}{@{\extracolsep{\fill}} l*{7}{c}} 
\hline
\hline
&\multicolumn{6}{c}{Dependent Variable: Culture-Talk (Fairly Often, Often=1)}\\
&Model 1&Model 2&Model 3&Model 4&Model 5&Model 6\\
&(Music)&(Music)&(Music)&(Music)&(Music)&(Music)\\
\hline
Listen to Music (Very Often=1)  &0.94 (4.98)&0.91 (4.57)&0.93 (4.67)&0.91 (4.57)&0.92 (4.62)&0.91 (4.60)\\ \\
Favorite Music?&&&&&\\
Opera, Classical=1&0.70 (3.03)&&&&&         \\
Jazz, Blues=1&&0.26 (1.12)&&&&         \\
Rock, Pop=1&&&0.12 (0.63)&&&         \\
Folk, Country=1&&&&-0.07 (-0.34)&&         \\
Christian, Gospel=1&&&&&-0.20 (-0.90)&         \\
Punk, Heavy Metal=1&&&&&&-0.23 (-0.58)\\ \\
$\rho_{21}$&-0.13&    -0.10&    -0.11&    -0.09&    -0.10&    -0.09\\
$\sigma^{2}_{\rho_{21}}$      & 0.11&     0.12&     0.12&     0.12&     0.12&     0.12\\
$\rho_{31}$  & 0.06&     0.13&     0.13&     0.09&     0.05&     0.04\\
$\sigma^{2}_{\rho_{31}}$ &     0.05&     0.05&     0.04&     0.05&     0.05&     0.09\\
$\rho_{32}$&-0.34&     0.04&    -0.05&    -0.01&     0.08&     0.28\\
$\sigma^{2}_{\rho_{32}}$      & 0.11&     0.12&     0.10&     0.11&     0.11&     0.19\\
$\chi^{2}$ test ($ \rho_{21}=\rho_{31}=\rho_{32}=0$)&    10.62&     8.31&     9.97&     4.84&     2.81&     2.56\\
N&\multicolumn{6}{c}{1756}\\
\hline
\end{tabular*}
\caption{}
\label{tab6}
\end{table}
\end{landscape}

\begin{landscape}
\begin{table}[htb!] 
\begin{tabular*}{\hsize}{@{\extracolsep{\fill}} l*{5}{c}}
\hline
\hline 
&\multicolumn{4}{c}{Dependent Variable: Culture-Talk (Fairly Often, Often=1)}\\
\hline
&Model 1&Model 2&Model 3&Model 4\\
&(Movies)&(Books)&(Arts)&(Plays)\\
\hline
Go to the Movies? (Fairly Often, Very Often=1)&1.09 (6.67)&&&         \\
Read Books? (Very Often=1)&&1.19 (7.13)&&         \\ 
Go to Museums? (Fairly Often, Very Often=1)&&&1.40 (4.73)&         \\
Go to Plays? (Fairly Often, Very Often=1)&&&&0.75 (2.84)\\ \\
Favorite Music?&&&&\\
Opera, Classical=1&0.41 (1.70)& -0.10 (-0.44)       &  0.59 (2.45)       & 0.56 (2.30)        \\ \\
\hline
&\multicolumn{4}{c}{Dependent Variable: Favorite Music? (Opera, Classical=1)}\\
\hline
&Model 1&Model 2&Model 3&Model 4\\
\hline
Go to the Movies? (Fairly Often, Very Often=1)&-0.11 (-0.51)&&&         \\
Read Books? (Very Often=1)&&0.40 (1.82)&&         \\
Go to Museums? (Fairly Often, Very Often=1)&&&0.36 (1.56)&         \\
Go to Plays? (Fairly Often, Very Often=1)&&&&0.44 (2.03)\\ \\
$\rho_{21}$ &    -0.20&    -0.25&    -0.15&    -0.07\\
$\sigma^{2}_{\rho_{21}}$      & 0.09&     0.10&     0.16&     0.15\\
$\rho_{31}$&-0.04&    -0.15&     0.03&    -0.13\\
$\sigma^{2}_{\rho_{31}}$      & 0.11&     0.13&     0.11&     0.11\\
$\rho_{32}$&-0.10&     0.17&    -0.20&    -0.19\\
$\sigma^{2}_{\rho_{32}}$&     0.12&     0.12&     0.10&     0.11\\
$\chi^{2}$ test ($ \rho_{21}=\rho_{31}=\rho_{32}=0$)&     6.49&     8.21&     5.22&     4.62\\
N&\multicolumn{4}{c}{1756}\\
\hline
\end{tabular*}
\caption{}
\label{tab7}
\end{table}
\end{landscape}



\newpage
\begin{figure}[here]
\centering
\subfigure[Univariate distribution of culture-mediated interaction probabilities (whole sample)]{
\includegraphics[width=6 in, height=4 in]{univariate-frequency-talk}
\label{fig1a}
}
\subfigure[Univariate distribution of culture-mediated interaction probabilities (culturally active subsample)]{         
 \includegraphics[width=6 in, height=4 in]{univariate-dist-hi-consump}
\label{fig1b}
}     
\caption{Univariate distribution of culture mediated interaction, 2002 and 2004 National Social Survey, Arts and Culture Supplement}
\end{figure}   

\newpage
\begin{figure}[here]
\centering
\subfigure[Predicted probabilities of culture-mediated interaction, by education]{
\includegraphics[width=6 in, height=4 in]{educ-pred-probs-interact-binary}
\label{fig2a}
}
\subfigure[Predicted probabilities of aesthetic consumption, by education]{         
\includegraphics[width=6 in, height=4 in]{educ-pred-probs-consump-binary}
\label{fig2b}
}     
\caption{Predicted probabilities of discussing cultural matters with friends and engaging in culture consumption for five cultural domains, 2002 and 2004 National Social Survey, Arts and Culture Supplement}
\end{figure}

\newpage
\begin{figure}[here]
\centering
\subfigure[Path model corresponding to the Seemingly Unrelated Bivariate Probit Model, with $Y_{1}$ and $Y_{2}$ are endogenous variables.]{
\includegraphics[width=2.8 in, height=2.1 in]{sur}
\label{fig3a}
}
\subfigure[\scriptsize Four models conceptualizing the connection between social position markers, aesthetic consumption and culture-mediated interaction.  In the figure $x_{1}$ and $x_{2}$ are vectors of socio-demographic characteristics (age, education, gender, ethnicity, etc.) and art consumption predictors (e.g. barriers to participation indicators, importance of art in the respondent's life, etc.), $\beta_{1}$ and $\beta_{2}$ are vectors of associated regression weights, $Y_{1}$ is a variable measuring frequency of culture mediated interaction in a given domain, $Y_{2}$ is a measure of aesthetic consumption in that domain and \emph{e} and \emph{v} represent all of the exogenous factors not correlated with socio-demographic characteristics over which both culture-mediated interaction and aesthetic consumption rates vary, respectively.  Finally, $\rho_{12}$ represents the correlation between exogenous variation in culture-mediated interaction and exogenous variation in aesthetic consumption beyond that which can be accounted for by the variables included in the $\beta X$ vector and the direct effect of interaction on consumption ($\delta_{12}$) and vice versa ($\delta_{21}$).]{         
 \includegraphics[width=6 in, height=4.5 in]{two-models}
\label{fig3b}}     
\caption{}
\end{figure}   




\newpage
\begin{figure}[here]
\centering
\includegraphics[width=6.5 in, height=7.5 in]{marg-effects-recursive}  
\caption{}
\label{fig4}
\end{figure}  

\newpage
\begin{figure}[here]
\centering
\includegraphics[width=6.5 in, height=4.5 in]{trivariate}    
\caption{}
\label{fig5}
\end{figure}    
 
 
 
\end{document}
